%==============================================================================
% Certificat de primalité inconditionnelle — premier de 29 998 chiffres
% Famille paramétrique  p = 3m(m+1)+1,  m = 2^a·3^b − 1
% Vérification numérique indépendante — mai 2026
%==============================================================================
%
% Utilisation dans le document principal :
%   %==============================================================================
% Certificat de primalité inconditionnelle — premier de 29 998 chiffres
% Famille paramétrique  p = 3m(m+1)+1,  m = 2^a·3^b − 1
% Vérification numérique indépendante — mai 2026
%==============================================================================
%
% Utilisation dans le document principal :
%   %==============================================================================
% Certificat de primalité inconditionnelle — premier de 29 998 chiffres
% Famille paramétrique  p = 3m(m+1)+1,  m = 2^a·3^b − 1
% Vérification numérique indépendante — mai 2026
%==============================================================================
%
% Utilisation dans le document principal :
%   %==============================================================================
% Certificat de primalité inconditionnelle — premier de 29 998 chiffres
% Famille paramétrique  p = 3m(m+1)+1,  m = 2^a·3^b − 1
% Vérification numérique indépendante — mai 2026
%==============================================================================
%
% Utilisation dans le document principal :
%   \input{certificat_29998_fr}
%
% Compilation autonome :
%   pdflatex certificat_29998_fr.tex
%
%==============================================================================

%------------------------------------------------------------------------------
\section{Certificat inconditionnel pour le premier de 29\,998 chiffres}
\label{sec:cert29998}
%------------------------------------------------------------------------------

Cette section présente le certificat de primalité Pocklington–Lehmer
complet et vérifié indépendamment pour le plus grand premier
actuellement produit par \textsc{PrimeQuest}.

%------------------------------------------------------------------------------
\subsection{Paramètres et structure}
%------------------------------------------------------------------------------

\begin{definition}[Instance de 29\,998 chiffres]\label{def:instance29998}
On pose
\[
  a = 19\,435,\qquad b = 19\,173,
\]
et on définit
\[
  m \;=\; 2^{19435}\cdot 3^{19173} - 1
    \qquad\bigl(14\,999 \text{ chiffres décimaux}\bigr),
\]
\[
  p \;=\; 3m(m+1)+1
    \qquad\bigl(29\,998 \text{ chiffres décimaux}\bigr).
\]
\end{definition}

Le rapport $b/a = 19\,173/19\,435 \approx 0{,}9865$ est voisin de l'optimum
théorique $\rho^* = \log_{10}2/\log_{10}3 \approx 0{,}6309$ minimisant le
nombre de chiffres à $a$ fixé ; la valeur retenue ici cible exactement
$29\,998$ chiffres et a été localisée par l'algorithme \textsc{PrimeQuest v4}
(voir Section~\ref{sec:algo}).

\medskip

\noindent\textbf{Identification.}\quad
Les premiers et derniers chiffres décimaux de $p$ sont :
\begin{equation}\label{eq:p29998chiffres}
  p \;=\; 1602087359509060348500037851549319056715\cdots
         16930852697403293697.
\end{equation}

%------------------------------------------------------------------------------
\subsection{Factorisation de $p-1$}
%------------------------------------------------------------------------------

\begin{rigorousbox}
\begin{lemma}[Factorisation de $p-1$]\label{lem:factorisation29998}
Avec $a, b, m, p$ comme dans la Définition~\ref{def:instance29998},
\[
  p - 1 \;=\; \underbrace{2^{19435}\cdot 3^{19174}}_{F}
              \;\cdot\; m,
\]
où $F := 2^{19435}\cdot 3^{19174}$ possède $15\,000$ chiffres décimaux et
vérifie $F > \sqrt{p}$.
\end{lemma}
\end{rigorousbox}

\begin{proof}
Puisque $m + 1 = 2^{19435}\cdot 3^{19173}$,
\[
  p - 1 \;=\; 3m(m+1) \;=\; 3m\cdot 2^{19435}\cdot 3^{19173}
           \;=\; 2^{19435}\cdot 3^{19174}\cdot m.
\]
Pour la borne $F > \sqrt{p}$, on note que $p < 3(m+1)^2$, donc
$\sqrt{p} < \sqrt{3}\,(m+1) < 3(m+1) = 2^{19435}\cdot 3^{19174} = F$.

\noindent Numériquement, en comparant les logarithmes en base $2$ :
\[
  2\log_2 F
  \;=\; 2\bigl(19435 + 19174\log_2 3\bigr)
  \;\approx\; 99\,650{,}14,
\]
\[
  \log_2 p \;\approx\; 99\,648{,}56,
\]
ce qui confirme $F^2 > p$ avec une marge de $\approx 1{,}58$ bit.
\end{proof}

%------------------------------------------------------------------------------
\subsection{Le certificat Pocklington–Lehmer}
%------------------------------------------------------------------------------

Le Théorème~\ref{thm:Pocklington} (rappelé Section~\ref{sec:family})
requiert, pour chaque premier $q \mid F$ — ici $q \in \{2, 3\}$ —, un entier
$w_q$ vérifiant
\[
  w_q^{p-1} \equiv 1 \pmod{p}
  \qquad\text{et}\qquad
  \gcd\!\bigl(w_q^{(p-1)/q} - 1,\; p\bigr) = 1.
\]

\begin{resultbox}[title={\textbf{Certificat} — premier de 29\,998 chiffres}]
\[
  p \;=\; 3m(m+1)+1,\qquad m = 2^{19435}\cdot 3^{19173}-1.
\]

\medskip
\noindent\textbf{Partie connue de $p-1$ :}
$\;F = 2^{19435}\cdot 3^{19174},\quad F > \sqrt{p}.\quad$\checkmark

\medskip
\renewcommand{\arraystretch}{1.4}
\begin{tabular}{@{}c l l@{}}
\toprule
Premier $q$ & Témoin $w_q$ & Conditions vérifiées \\
\midrule
$2$ & $w_2 = 5$ &
  $5^{p-1}\equiv 1\pmod{p}$\quad et\quad
  $\gcd(5^{(p-1)/2}-1,\,p)=1$ \\
$3$ & $w_3 = 7$ &
  $7^{p-1}\equiv 1\pmod{p}$\quad et\quad
  $\gcd(7^{(p-1)/3}-1,\,p)=1$ \\
\bottomrule
\end{tabular}

\medskip
\noindent\textbf{Conclusion :} $p$ est premier (preuve déterministe). \checkmark
\end{resultbox}

%------------------------------------------------------------------------------
\subsection{Vérification numérique indépendante}
%------------------------------------------------------------------------------

Les quatre conditions du certificat ont été vérifiées de façon indépendante
à l'aide de \texttt{Python 3.12} et de son arithmétique entière en précision
arbitraire (backend GMP). Le script ne comporte aucune étape probabiliste ;
toutes les opérations sont des arithmétiques modulaires exactes.

\begin{rigorousbox}[title={\textbf{Attestation de vérification} — calcul indépendant}]
\textbf{Environnement :} Python~3.12,
\texttt{sys.set\_int\_max\_str\_digits(40000)} ;
exponentiations modulaires via \texttt{pow(base, exp, mod)}.\\[4pt]

\textbf{Étape 1 — Reconstruction de $p$.}\quad
Calcul de $m = 2^{19435}\cdot 3^{19173}-1$ (14\,999 chiffres) et de
$p = 3m(m+1)+1$ (29\,998 chiffres) à partir des paramètres bruts.
Premiers 40 chiffres de $p$ :
\texttt{1602087359509060348500037851549319056715} ;
20 derniers : \texttt{16930852697403293697}.
Correspond exactement au certificat stocké.\\[4pt]

\textbf{Étape 2 — Borne Pocklington.}
\[
  2\log_2 F = 99\,650{,}1420 \;>\; \log_2 p = 99\,648{,}5570.\quad\checkmark
\]

\textbf{Étape 3 — Témoin $w_2 = 5$ pour $q = 2$.}
\begin{align*}
  5^{p-1} &\equiv 1 \pmod{p},\quad\checkmark \\
  \gcd(5^{(p-1)/2}-1,\,p) &= 1.\quad\checkmark
\end{align*}

\textbf{Étape 4 — Témoin $w_3 = 7$ pour $q = 3$.}
\begin{align*}
  7^{p-1} &\equiv 1 \pmod{p},\quad\checkmark \\
  \gcd(7^{(p-1)/3}-1,\,p) &= 1.\quad\checkmark
\end{align*}

\textbf{Durée totale :} $6\,945$\,s $\approx 1$\,h\,$56$\,min
(cœur unique, Python GMP, sans parallélisme).\\[4pt]

\textbf{Conclusion :} les quatre conditions Pocklington–Lehmer sont satisfaites.
Le nombre $p$ est \emph{premier de façon inconditionnelle}.
\end{rigorousbox}

%------------------------------------------------------------------------------
\subsection{Comparaison avec le calcul original}
%------------------------------------------------------------------------------

\begin{table}[h]
\centering
\caption{Recherche originale vs.\ vérification indépendante du certificat
         pour le premier de 29\,998 chiffres.}
\label{tab:comparaison29998}
\renewcommand{\arraystretch}{1.3}
\begin{tabular}{@{}lll@{}}
\toprule
 & \textbf{PrimeQuest v4 (recherche)} & \textbf{Vérification indépendante} \\
\midrule
Objectif       & Trouver un candidat premier         & Vérifier le certificat \\
Machine        & ARM64, 9 cœurs (Snapdragon X)       & Cœur unique (Python GMP) \\
Durée          & $11\,146$\,s ($3$\,h\,$06$\,min)   & $6\,945$\,s ($1$\,h\,$56$\,min) \\
Tests MR       & $286$ (pré-filtre probabiliste)     & — (aucune étape probabiliste) \\
Résultat       & Certificat produit                  & Certificat confirmé \\
\bottomrule
\end{tabular}
\end{table}

La vérification est $\approx 1{,}6$ fois plus rapide que la recherche
originale car elle exécute exactement quatre exponentiations modulaires sur
le candidat connu, sans crible, sans filtre du Théorème~2 ni
pré-sélection Miller–Rabin.

%------------------------------------------------------------------------------
\subsection{Rang dans les records}
%------------------------------------------------------------------------------

Le Tableau~\ref{tab:tousresultats} (Section~\ref{sec:resultats}) liste
l'ensemble des quatre certificats \textsc{PrimeQuest}. L'instance de
$29\,998$ chiffres est actuellement le plus grand certificat
Pocklington–Lehmer inconditionnel dans la famille
$p = 3m(m+1)+1$, $m = 2^a\cdot 3^b - 1$,
obtenu sur un ordinateur portable grand public sans calcul distribué ni
matériel spécialisé.

%==============================================================================
% PRÉAMBULE MINIMAL AUTONOME (décommenter pour compilation stand-alone)
%==============================================================================
% \iffalse   % <-- supprimer cette ligne pour activer le mode autonome
% \documentclass[11pt,a4paper]{article}
% \usepackage[utf8]{inputenc}
% \usepackage[T1]{fontenc}
% \usepackage[french]{babel}
% \usepackage{amsmath,amssymb,amsthm,mathtools,booktabs}
% \usepackage{geometry}\geometry{margin=2.5cm}
% \usepackage[dvipsnames]{xcolor}
% \usepackage{tcolorbox}
% \tcbuselibrary{skins,breakable}
% \newtcolorbox{rigorousbox}[1][]{colback=Green!8,colframe=Green!60!black,
%   arc=1mm,title=\textbf{Rigoureux},fonttitle=\bfseries,breakable,#1}
% \newtcolorbox{resultbox}[1][]{colback=RoyalBlue!5,colframe=RoyalBlue!70!black,
%   arc=1mm,title=\textbf{Résultat calculatoire},fonttitle=\bfseries,breakable,#1}
% \newtheorem{theorem}{Théorème}[section]
% \newtheorem{lemma}[theorem]{Lemme}
% \theoremstyle{definition}
% \newtheorem{definition}[theorem]{Définition}
% \begin{document}
% \input{certificat_29998_fr}
% \end{document}
% \fi

%
% Compilation autonome :
%   pdflatex certificat_29998_fr.tex
%
%==============================================================================

%------------------------------------------------------------------------------
\section{Certificat inconditionnel pour le premier de 29\,998 chiffres}
\label{sec:cert29998}
%------------------------------------------------------------------------------

Cette section présente le certificat de primalité Pocklington–Lehmer
complet et vérifié indépendamment pour le plus grand premier
actuellement produit par \textsc{PrimeQuest}.

%------------------------------------------------------------------------------
\subsection{Paramètres et structure}
%------------------------------------------------------------------------------

\begin{definition}[Instance de 29\,998 chiffres]\label{def:instance29998}
On pose
\[
  a = 19\,435,\qquad b = 19\,173,
\]
et on définit
\[
  m \;=\; 2^{19435}\cdot 3^{19173} - 1
    \qquad\bigl(14\,999 \text{ chiffres décimaux}\bigr),
\]
\[
  p \;=\; 3m(m+1)+1
    \qquad\bigl(29\,998 \text{ chiffres décimaux}\bigr).
\]
\end{definition}

Le rapport $b/a = 19\,173/19\,435 \approx 0{,}9865$ est voisin de l'optimum
théorique $\rho^* = \log_{10}2/\log_{10}3 \approx 0{,}6309$ minimisant le
nombre de chiffres à $a$ fixé ; la valeur retenue ici cible exactement
$29\,998$ chiffres et a été localisée par l'algorithme \textsc{PrimeQuest v4}
(voir Section~\ref{sec:algo}).

\medskip

\noindent\textbf{Identification.}\quad
Les premiers et derniers chiffres décimaux de $p$ sont :
\begin{equation}\label{eq:p29998chiffres}
  p \;=\; 1602087359509060348500037851549319056715\cdots
         16930852697403293697.
\end{equation}

%------------------------------------------------------------------------------
\subsection{Factorisation de $p-1$}
%------------------------------------------------------------------------------

\begin{rigorousbox}
\begin{lemma}[Factorisation de $p-1$]\label{lem:factorisation29998}
Avec $a, b, m, p$ comme dans la Définition~\ref{def:instance29998},
\[
  p - 1 \;=\; \underbrace{2^{19435}\cdot 3^{19174}}_{F}
              \;\cdot\; m,
\]
où $F := 2^{19435}\cdot 3^{19174}$ possède $15\,000$ chiffres décimaux et
vérifie $F > \sqrt{p}$.
\end{lemma}
\end{rigorousbox}

\begin{proof}
Puisque $m + 1 = 2^{19435}\cdot 3^{19173}$,
\[
  p - 1 \;=\; 3m(m+1) \;=\; 3m\cdot 2^{19435}\cdot 3^{19173}
           \;=\; 2^{19435}\cdot 3^{19174}\cdot m.
\]
Pour la borne $F > \sqrt{p}$, on note que $p < 3(m+1)^2$, donc
$\sqrt{p} < \sqrt{3}\,(m+1) < 3(m+1) = 2^{19435}\cdot 3^{19174} = F$.

\noindent Numériquement, en comparant les logarithmes en base $2$ :
\[
  2\log_2 F
  \;=\; 2\bigl(19435 + 19174\log_2 3\bigr)
  \;\approx\; 99\,650{,}14,
\]
\[
  \log_2 p \;\approx\; 99\,648{,}56,
\]
ce qui confirme $F^2 > p$ avec une marge de $\approx 1{,}58$ bit.
\end{proof}

%------------------------------------------------------------------------------
\subsection{Le certificat Pocklington–Lehmer}
%------------------------------------------------------------------------------

Le Théorème~\ref{thm:Pocklington} (rappelé Section~\ref{sec:family})
requiert, pour chaque premier $q \mid F$ — ici $q \in \{2, 3\}$ —, un entier
$w_q$ vérifiant
\[
  w_q^{p-1} \equiv 1 \pmod{p}
  \qquad\text{et}\qquad
  \gcd\!\bigl(w_q^{(p-1)/q} - 1,\; p\bigr) = 1.
\]

\begin{resultbox}[title={\textbf{Certificat} — premier de 29\,998 chiffres}]
\[
  p \;=\; 3m(m+1)+1,\qquad m = 2^{19435}\cdot 3^{19173}-1.
\]

\medskip
\noindent\textbf{Partie connue de $p-1$ :}
$\;F = 2^{19435}\cdot 3^{19174},\quad F > \sqrt{p}.\quad$\checkmark

\medskip
\renewcommand{\arraystretch}{1.4}
\begin{tabular}{@{}c l l@{}}
\toprule
Premier $q$ & Témoin $w_q$ & Conditions vérifiées \\
\midrule
$2$ & $w_2 = 5$ &
  $5^{p-1}\equiv 1\pmod{p}$\quad et\quad
  $\gcd(5^{(p-1)/2}-1,\,p)=1$ \\
$3$ & $w_3 = 7$ &
  $7^{p-1}\equiv 1\pmod{p}$\quad et\quad
  $\gcd(7^{(p-1)/3}-1,\,p)=1$ \\
\bottomrule
\end{tabular}

\medskip
\noindent\textbf{Conclusion :} $p$ est premier (preuve déterministe). \checkmark
\end{resultbox}

%------------------------------------------------------------------------------
\subsection{Vérification numérique indépendante}
%------------------------------------------------------------------------------

Les quatre conditions du certificat ont été vérifiées de façon indépendante
à l'aide de \texttt{Python 3.12} et de son arithmétique entière en précision
arbitraire (backend GMP). Le script ne comporte aucune étape probabiliste ;
toutes les opérations sont des arithmétiques modulaires exactes.

\begin{rigorousbox}[title={\textbf{Attestation de vérification} — calcul indépendant}]
\textbf{Environnement :} Python~3.12,
\texttt{sys.set\_int\_max\_str\_digits(40000)} ;
exponentiations modulaires via \texttt{pow(base, exp, mod)}.\\[4pt]

\textbf{Étape 1 — Reconstruction de $p$.}\quad
Calcul de $m = 2^{19435}\cdot 3^{19173}-1$ (14\,999 chiffres) et de
$p = 3m(m+1)+1$ (29\,998 chiffres) à partir des paramètres bruts.
Premiers 40 chiffres de $p$ :
\texttt{1602087359509060348500037851549319056715} ;
20 derniers : \texttt{16930852697403293697}.
Correspond exactement au certificat stocké.\\[4pt]

\textbf{Étape 2 — Borne Pocklington.}
\[
  2\log_2 F = 99\,650{,}1420 \;>\; \log_2 p = 99\,648{,}5570.\quad\checkmark
\]

\textbf{Étape 3 — Témoin $w_2 = 5$ pour $q = 2$.}
\begin{align*}
  5^{p-1} &\equiv 1 \pmod{p},\quad\checkmark \\
  \gcd(5^{(p-1)/2}-1,\,p) &= 1.\quad\checkmark
\end{align*}

\textbf{Étape 4 — Témoin $w_3 = 7$ pour $q = 3$.}
\begin{align*}
  7^{p-1} &\equiv 1 \pmod{p},\quad\checkmark \\
  \gcd(7^{(p-1)/3}-1,\,p) &= 1.\quad\checkmark
\end{align*}

\textbf{Durée totale :} $6\,945$\,s $\approx 1$\,h\,$56$\,min
(cœur unique, Python GMP, sans parallélisme).\\[4pt]

\textbf{Conclusion :} les quatre conditions Pocklington–Lehmer sont satisfaites.
Le nombre $p$ est \emph{premier de façon inconditionnelle}.
\end{rigorousbox}

%------------------------------------------------------------------------------
\subsection{Comparaison avec le calcul original}
%------------------------------------------------------------------------------

\begin{table}[h]
\centering
\caption{Recherche originale vs.\ vérification indépendante du certificat
         pour le premier de 29\,998 chiffres.}
\label{tab:comparaison29998}
\renewcommand{\arraystretch}{1.3}
\begin{tabular}{@{}lll@{}}
\toprule
 & \textbf{PrimeQuest v4 (recherche)} & \textbf{Vérification indépendante} \\
\midrule
Objectif       & Trouver un candidat premier         & Vérifier le certificat \\
Machine        & ARM64, 9 cœurs (Snapdragon X)       & Cœur unique (Python GMP) \\
Durée          & $11\,146$\,s ($3$\,h\,$06$\,min)   & $6\,945$\,s ($1$\,h\,$56$\,min) \\
Tests MR       & $286$ (pré-filtre probabiliste)     & — (aucune étape probabiliste) \\
Résultat       & Certificat produit                  & Certificat confirmé \\
\bottomrule
\end{tabular}
\end{table}

La vérification est $\approx 1{,}6$ fois plus rapide que la recherche
originale car elle exécute exactement quatre exponentiations modulaires sur
le candidat connu, sans crible, sans filtre du Théorème~2 ni
pré-sélection Miller–Rabin.

%------------------------------------------------------------------------------
\subsection{Rang dans les records}
%------------------------------------------------------------------------------

Le Tableau~\ref{tab:tousresultats} (Section~\ref{sec:resultats}) liste
l'ensemble des quatre certificats \textsc{PrimeQuest}. L'instance de
$29\,998$ chiffres est actuellement le plus grand certificat
Pocklington–Lehmer inconditionnel dans la famille
$p = 3m(m+1)+1$, $m = 2^a\cdot 3^b - 1$,
obtenu sur un ordinateur portable grand public sans calcul distribué ni
matériel spécialisé.

%==============================================================================
% PRÉAMBULE MINIMAL AUTONOME (décommenter pour compilation stand-alone)
%==============================================================================
% \iffalse   % <-- supprimer cette ligne pour activer le mode autonome
% \documentclass[11pt,a4paper]{article}
% \usepackage[utf8]{inputenc}
% \usepackage[T1]{fontenc}
% \usepackage[french]{babel}
% \usepackage{amsmath,amssymb,amsthm,mathtools,booktabs}
% \usepackage{geometry}\geometry{margin=2.5cm}
% \usepackage[dvipsnames]{xcolor}
% \usepackage{tcolorbox}
% \tcbuselibrary{skins,breakable}
% \newtcolorbox{rigorousbox}[1][]{colback=Green!8,colframe=Green!60!black,
%   arc=1mm,title=\textbf{Rigoureux},fonttitle=\bfseries,breakable,#1}
% \newtcolorbox{resultbox}[1][]{colback=RoyalBlue!5,colframe=RoyalBlue!70!black,
%   arc=1mm,title=\textbf{Résultat calculatoire},fonttitle=\bfseries,breakable,#1}
% \newtheorem{theorem}{Théorème}[section]
% \newtheorem{lemma}[theorem]{Lemme}
% \theoremstyle{definition}
% \newtheorem{definition}[theorem]{Définition}
% \begin{document}
% %==============================================================================
% Certificat de primalité inconditionnelle — premier de 29 998 chiffres
% Famille paramétrique  p = 3m(m+1)+1,  m = 2^a·3^b − 1
% Vérification numérique indépendante — mai 2026
%==============================================================================
%
% Utilisation dans le document principal :
%   \input{certificat_29998_fr}
%
% Compilation autonome :
%   pdflatex certificat_29998_fr.tex
%
%==============================================================================

%------------------------------------------------------------------------------
\section{Certificat inconditionnel pour le premier de 29\,998 chiffres}
\label{sec:cert29998}
%------------------------------------------------------------------------------

Cette section présente le certificat de primalité Pocklington–Lehmer
complet et vérifié indépendamment pour le plus grand premier
actuellement produit par \textsc{PrimeQuest}.

%------------------------------------------------------------------------------
\subsection{Paramètres et structure}
%------------------------------------------------------------------------------

\begin{definition}[Instance de 29\,998 chiffres]\label{def:instance29998}
On pose
\[
  a = 19\,435,\qquad b = 19\,173,
\]
et on définit
\[
  m \;=\; 2^{19435}\cdot 3^{19173} - 1
    \qquad\bigl(14\,999 \text{ chiffres décimaux}\bigr),
\]
\[
  p \;=\; 3m(m+1)+1
    \qquad\bigl(29\,998 \text{ chiffres décimaux}\bigr).
\]
\end{definition}

Le rapport $b/a = 19\,173/19\,435 \approx 0{,}9865$ est voisin de l'optimum
théorique $\rho^* = \log_{10}2/\log_{10}3 \approx 0{,}6309$ minimisant le
nombre de chiffres à $a$ fixé ; la valeur retenue ici cible exactement
$29\,998$ chiffres et a été localisée par l'algorithme \textsc{PrimeQuest v4}
(voir Section~\ref{sec:algo}).

\medskip

\noindent\textbf{Identification.}\quad
Les premiers et derniers chiffres décimaux de $p$ sont :
\begin{equation}\label{eq:p29998chiffres}
  p \;=\; 1602087359509060348500037851549319056715\cdots
         16930852697403293697.
\end{equation}

%------------------------------------------------------------------------------
\subsection{Factorisation de $p-1$}
%------------------------------------------------------------------------------

\begin{rigorousbox}
\begin{lemma}[Factorisation de $p-1$]\label{lem:factorisation29998}
Avec $a, b, m, p$ comme dans la Définition~\ref{def:instance29998},
\[
  p - 1 \;=\; \underbrace{2^{19435}\cdot 3^{19174}}_{F}
              \;\cdot\; m,
\]
où $F := 2^{19435}\cdot 3^{19174}$ possède $15\,000$ chiffres décimaux et
vérifie $F > \sqrt{p}$.
\end{lemma}
\end{rigorousbox}

\begin{proof}
Puisque $m + 1 = 2^{19435}\cdot 3^{19173}$,
\[
  p - 1 \;=\; 3m(m+1) \;=\; 3m\cdot 2^{19435}\cdot 3^{19173}
           \;=\; 2^{19435}\cdot 3^{19174}\cdot m.
\]
Pour la borne $F > \sqrt{p}$, on note que $p < 3(m+1)^2$, donc
$\sqrt{p} < \sqrt{3}\,(m+1) < 3(m+1) = 2^{19435}\cdot 3^{19174} = F$.

\noindent Numériquement, en comparant les logarithmes en base $2$ :
\[
  2\log_2 F
  \;=\; 2\bigl(19435 + 19174\log_2 3\bigr)
  \;\approx\; 99\,650{,}14,
\]
\[
  \log_2 p \;\approx\; 99\,648{,}56,
\]
ce qui confirme $F^2 > p$ avec une marge de $\approx 1{,}58$ bit.
\end{proof}

%------------------------------------------------------------------------------
\subsection{Le certificat Pocklington–Lehmer}
%------------------------------------------------------------------------------

Le Théorème~\ref{thm:Pocklington} (rappelé Section~\ref{sec:family})
requiert, pour chaque premier $q \mid F$ — ici $q \in \{2, 3\}$ —, un entier
$w_q$ vérifiant
\[
  w_q^{p-1} \equiv 1 \pmod{p}
  \qquad\text{et}\qquad
  \gcd\!\bigl(w_q^{(p-1)/q} - 1,\; p\bigr) = 1.
\]

\begin{resultbox}[title={\textbf{Certificat} — premier de 29\,998 chiffres}]
\[
  p \;=\; 3m(m+1)+1,\qquad m = 2^{19435}\cdot 3^{19173}-1.
\]

\medskip
\noindent\textbf{Partie connue de $p-1$ :}
$\;F = 2^{19435}\cdot 3^{19174},\quad F > \sqrt{p}.\quad$\checkmark

\medskip
\renewcommand{\arraystretch}{1.4}
\begin{tabular}{@{}c l l@{}}
\toprule
Premier $q$ & Témoin $w_q$ & Conditions vérifiées \\
\midrule
$2$ & $w_2 = 5$ &
  $5^{p-1}\equiv 1\pmod{p}$\quad et\quad
  $\gcd(5^{(p-1)/2}-1,\,p)=1$ \\
$3$ & $w_3 = 7$ &
  $7^{p-1}\equiv 1\pmod{p}$\quad et\quad
  $\gcd(7^{(p-1)/3}-1,\,p)=1$ \\
\bottomrule
\end{tabular}

\medskip
\noindent\textbf{Conclusion :} $p$ est premier (preuve déterministe). \checkmark
\end{resultbox}

%------------------------------------------------------------------------------
\subsection{Vérification numérique indépendante}
%------------------------------------------------------------------------------

Les quatre conditions du certificat ont été vérifiées de façon indépendante
à l'aide de \texttt{Python 3.12} et de son arithmétique entière en précision
arbitraire (backend GMP). Le script ne comporte aucune étape probabiliste ;
toutes les opérations sont des arithmétiques modulaires exactes.

\begin{rigorousbox}[title={\textbf{Attestation de vérification} — calcul indépendant}]
\textbf{Environnement :} Python~3.12,
\texttt{sys.set\_int\_max\_str\_digits(40000)} ;
exponentiations modulaires via \texttt{pow(base, exp, mod)}.\\[4pt]

\textbf{Étape 1 — Reconstruction de $p$.}\quad
Calcul de $m = 2^{19435}\cdot 3^{19173}-1$ (14\,999 chiffres) et de
$p = 3m(m+1)+1$ (29\,998 chiffres) à partir des paramètres bruts.
Premiers 40 chiffres de $p$ :
\texttt{1602087359509060348500037851549319056715} ;
20 derniers : \texttt{16930852697403293697}.
Correspond exactement au certificat stocké.\\[4pt]

\textbf{Étape 2 — Borne Pocklington.}
\[
  2\log_2 F = 99\,650{,}1420 \;>\; \log_2 p = 99\,648{,}5570.\quad\checkmark
\]

\textbf{Étape 3 — Témoin $w_2 = 5$ pour $q = 2$.}
\begin{align*}
  5^{p-1} &\equiv 1 \pmod{p},\quad\checkmark \\
  \gcd(5^{(p-1)/2}-1,\,p) &= 1.\quad\checkmark
\end{align*}

\textbf{Étape 4 — Témoin $w_3 = 7$ pour $q = 3$.}
\begin{align*}
  7^{p-1} &\equiv 1 \pmod{p},\quad\checkmark \\
  \gcd(7^{(p-1)/3}-1,\,p) &= 1.\quad\checkmark
\end{align*}

\textbf{Durée totale :} $6\,945$\,s $\approx 1$\,h\,$56$\,min
(cœur unique, Python GMP, sans parallélisme).\\[4pt]

\textbf{Conclusion :} les quatre conditions Pocklington–Lehmer sont satisfaites.
Le nombre $p$ est \emph{premier de façon inconditionnelle}.
\end{rigorousbox}

%------------------------------------------------------------------------------
\subsection{Comparaison avec le calcul original}
%------------------------------------------------------------------------------

\begin{table}[h]
\centering
\caption{Recherche originale vs.\ vérification indépendante du certificat
         pour le premier de 29\,998 chiffres.}
\label{tab:comparaison29998}
\renewcommand{\arraystretch}{1.3}
\begin{tabular}{@{}lll@{}}
\toprule
 & \textbf{PrimeQuest v4 (recherche)} & \textbf{Vérification indépendante} \\
\midrule
Objectif       & Trouver un candidat premier         & Vérifier le certificat \\
Machine        & ARM64, 9 cœurs (Snapdragon X)       & Cœur unique (Python GMP) \\
Durée          & $11\,146$\,s ($3$\,h\,$06$\,min)   & $6\,945$\,s ($1$\,h\,$56$\,min) \\
Tests MR       & $286$ (pré-filtre probabiliste)     & — (aucune étape probabiliste) \\
Résultat       & Certificat produit                  & Certificat confirmé \\
\bottomrule
\end{tabular}
\end{table}

La vérification est $\approx 1{,}6$ fois plus rapide que la recherche
originale car elle exécute exactement quatre exponentiations modulaires sur
le candidat connu, sans crible, sans filtre du Théorème~2 ni
pré-sélection Miller–Rabin.

%------------------------------------------------------------------------------
\subsection{Rang dans les records}
%------------------------------------------------------------------------------

Le Tableau~\ref{tab:tousresultats} (Section~\ref{sec:resultats}) liste
l'ensemble des quatre certificats \textsc{PrimeQuest}. L'instance de
$29\,998$ chiffres est actuellement le plus grand certificat
Pocklington–Lehmer inconditionnel dans la famille
$p = 3m(m+1)+1$, $m = 2^a\cdot 3^b - 1$,
obtenu sur un ordinateur portable grand public sans calcul distribué ni
matériel spécialisé.

%==============================================================================
% PRÉAMBULE MINIMAL AUTONOME (décommenter pour compilation stand-alone)
%==============================================================================
% \iffalse   % <-- supprimer cette ligne pour activer le mode autonome
% \documentclass[11pt,a4paper]{article}
% \usepackage[utf8]{inputenc}
% \usepackage[T1]{fontenc}
% \usepackage[french]{babel}
% \usepackage{amsmath,amssymb,amsthm,mathtools,booktabs}
% \usepackage{geometry}\geometry{margin=2.5cm}
% \usepackage[dvipsnames]{xcolor}
% \usepackage{tcolorbox}
% \tcbuselibrary{skins,breakable}
% \newtcolorbox{rigorousbox}[1][]{colback=Green!8,colframe=Green!60!black,
%   arc=1mm,title=\textbf{Rigoureux},fonttitle=\bfseries,breakable,#1}
% \newtcolorbox{resultbox}[1][]{colback=RoyalBlue!5,colframe=RoyalBlue!70!black,
%   arc=1mm,title=\textbf{Résultat calculatoire},fonttitle=\bfseries,breakable,#1}
% \newtheorem{theorem}{Théorème}[section]
% \newtheorem{lemma}[theorem]{Lemme}
% \theoremstyle{definition}
% \newtheorem{definition}[theorem]{Définition}
% \begin{document}
% \input{certificat_29998_fr}
% \end{document}
% \fi

% \end{document}
% \fi

%
% Compilation autonome :
%   pdflatex certificat_29998_fr.tex
%
%==============================================================================

%------------------------------------------------------------------------------
\section{Certificat inconditionnel pour le premier de 29\,998 chiffres}
\label{sec:cert29998}
%------------------------------------------------------------------------------

Cette section présente le certificat de primalité Pocklington–Lehmer
complet et vérifié indépendamment pour le plus grand premier
actuellement produit par \textsc{PrimeQuest}.

%------------------------------------------------------------------------------
\subsection{Paramètres et structure}
%------------------------------------------------------------------------------

\begin{definition}[Instance de 29\,998 chiffres]\label{def:instance29998}
On pose
\[
  a = 19\,435,\qquad b = 19\,173,
\]
et on définit
\[
  m \;=\; 2^{19435}\cdot 3^{19173} - 1
    \qquad\bigl(14\,999 \text{ chiffres décimaux}\bigr),
\]
\[
  p \;=\; 3m(m+1)+1
    \qquad\bigl(29\,998 \text{ chiffres décimaux}\bigr).
\]
\end{definition}

Le rapport $b/a = 19\,173/19\,435 \approx 0{,}9865$ est voisin de l'optimum
théorique $\rho^* = \log_{10}2/\log_{10}3 \approx 0{,}6309$ minimisant le
nombre de chiffres à $a$ fixé ; la valeur retenue ici cible exactement
$29\,998$ chiffres et a été localisée par l'algorithme \textsc{PrimeQuest v4}
(voir Section~\ref{sec:algo}).

\medskip

\noindent\textbf{Identification.}\quad
Les premiers et derniers chiffres décimaux de $p$ sont :
\begin{equation}\label{eq:p29998chiffres}
  p \;=\; 1602087359509060348500037851549319056715\cdots
         16930852697403293697.
\end{equation}

%------------------------------------------------------------------------------
\subsection{Factorisation de $p-1$}
%------------------------------------------------------------------------------

\begin{rigorousbox}
\begin{lemma}[Factorisation de $p-1$]\label{lem:factorisation29998}
Avec $a, b, m, p$ comme dans la Définition~\ref{def:instance29998},
\[
  p - 1 \;=\; \underbrace{2^{19435}\cdot 3^{19174}}_{F}
              \;\cdot\; m,
\]
où $F := 2^{19435}\cdot 3^{19174}$ possède $15\,000$ chiffres décimaux et
vérifie $F > \sqrt{p}$.
\end{lemma}
\end{rigorousbox}

\begin{proof}
Puisque $m + 1 = 2^{19435}\cdot 3^{19173}$,
\[
  p - 1 \;=\; 3m(m+1) \;=\; 3m\cdot 2^{19435}\cdot 3^{19173}
           \;=\; 2^{19435}\cdot 3^{19174}\cdot m.
\]
Pour la borne $F > \sqrt{p}$, on note que $p < 3(m+1)^2$, donc
$\sqrt{p} < \sqrt{3}\,(m+1) < 3(m+1) = 2^{19435}\cdot 3^{19174} = F$.

\noindent Numériquement, en comparant les logarithmes en base $2$ :
\[
  2\log_2 F
  \;=\; 2\bigl(19435 + 19174\log_2 3\bigr)
  \;\approx\; 99\,650{,}14,
\]
\[
  \log_2 p \;\approx\; 99\,648{,}56,
\]
ce qui confirme $F^2 > p$ avec une marge de $\approx 1{,}58$ bit.
\end{proof}

%------------------------------------------------------------------------------
\subsection{Le certificat Pocklington–Lehmer}
%------------------------------------------------------------------------------

Le Théorème~\ref{thm:Pocklington} (rappelé Section~\ref{sec:family})
requiert, pour chaque premier $q \mid F$ — ici $q \in \{2, 3\}$ —, un entier
$w_q$ vérifiant
\[
  w_q^{p-1} \equiv 1 \pmod{p}
  \qquad\text{et}\qquad
  \gcd\!\bigl(w_q^{(p-1)/q} - 1,\; p\bigr) = 1.
\]

\begin{resultbox}[title={\textbf{Certificat} — premier de 29\,998 chiffres}]
\[
  p \;=\; 3m(m+1)+1,\qquad m = 2^{19435}\cdot 3^{19173}-1.
\]

\medskip
\noindent\textbf{Partie connue de $p-1$ :}
$\;F = 2^{19435}\cdot 3^{19174},\quad F > \sqrt{p}.\quad$\checkmark

\medskip
\renewcommand{\arraystretch}{1.4}
\begin{tabular}{@{}c l l@{}}
\toprule
Premier $q$ & Témoin $w_q$ & Conditions vérifiées \\
\midrule
$2$ & $w_2 = 5$ &
  $5^{p-1}\equiv 1\pmod{p}$\quad et\quad
  $\gcd(5^{(p-1)/2}-1,\,p)=1$ \\
$3$ & $w_3 = 7$ &
  $7^{p-1}\equiv 1\pmod{p}$\quad et\quad
  $\gcd(7^{(p-1)/3}-1,\,p)=1$ \\
\bottomrule
\end{tabular}

\medskip
\noindent\textbf{Conclusion :} $p$ est premier (preuve déterministe). \checkmark
\end{resultbox}

%------------------------------------------------------------------------------
\subsection{Vérification numérique indépendante}
%------------------------------------------------------------------------------

Les quatre conditions du certificat ont été vérifiées de façon indépendante
à l'aide de \texttt{Python 3.12} et de son arithmétique entière en précision
arbitraire (backend GMP). Le script ne comporte aucune étape probabiliste ;
toutes les opérations sont des arithmétiques modulaires exactes.

\begin{rigorousbox}[title={\textbf{Attestation de vérification} — calcul indépendant}]
\textbf{Environnement :} Python~3.12,
\texttt{sys.set\_int\_max\_str\_digits(40000)} ;
exponentiations modulaires via \texttt{pow(base, exp, mod)}.\\[4pt]

\textbf{Étape 1 — Reconstruction de $p$.}\quad
Calcul de $m = 2^{19435}\cdot 3^{19173}-1$ (14\,999 chiffres) et de
$p = 3m(m+1)+1$ (29\,998 chiffres) à partir des paramètres bruts.
Premiers 40 chiffres de $p$ :
\texttt{1602087359509060348500037851549319056715} ;
20 derniers : \texttt{16930852697403293697}.
Correspond exactement au certificat stocké.\\[4pt]

\textbf{Étape 2 — Borne Pocklington.}
\[
  2\log_2 F = 99\,650{,}1420 \;>\; \log_2 p = 99\,648{,}5570.\quad\checkmark
\]

\textbf{Étape 3 — Témoin $w_2 = 5$ pour $q = 2$.}
\begin{align*}
  5^{p-1} &\equiv 1 \pmod{p},\quad\checkmark \\
  \gcd(5^{(p-1)/2}-1,\,p) &= 1.\quad\checkmark
\end{align*}

\textbf{Étape 4 — Témoin $w_3 = 7$ pour $q = 3$.}
\begin{align*}
  7^{p-1} &\equiv 1 \pmod{p},\quad\checkmark \\
  \gcd(7^{(p-1)/3}-1,\,p) &= 1.\quad\checkmark
\end{align*}

\textbf{Durée totale :} $6\,945$\,s $\approx 1$\,h\,$56$\,min
(cœur unique, Python GMP, sans parallélisme).\\[4pt]

\textbf{Conclusion :} les quatre conditions Pocklington–Lehmer sont satisfaites.
Le nombre $p$ est \emph{premier de façon inconditionnelle}.
\end{rigorousbox}

%------------------------------------------------------------------------------
\subsection{Comparaison avec le calcul original}
%------------------------------------------------------------------------------

\begin{table}[h]
\centering
\caption{Recherche originale vs.\ vérification indépendante du certificat
         pour le premier de 29\,998 chiffres.}
\label{tab:comparaison29998}
\renewcommand{\arraystretch}{1.3}
\begin{tabular}{@{}lll@{}}
\toprule
 & \textbf{PrimeQuest v4 (recherche)} & \textbf{Vérification indépendante} \\
\midrule
Objectif       & Trouver un candidat premier         & Vérifier le certificat \\
Machine        & ARM64, 9 cœurs (Snapdragon X)       & Cœur unique (Python GMP) \\
Durée          & $11\,146$\,s ($3$\,h\,$06$\,min)   & $6\,945$\,s ($1$\,h\,$56$\,min) \\
Tests MR       & $286$ (pré-filtre probabiliste)     & — (aucune étape probabiliste) \\
Résultat       & Certificat produit                  & Certificat confirmé \\
\bottomrule
\end{tabular}
\end{table}

La vérification est $\approx 1{,}6$ fois plus rapide que la recherche
originale car elle exécute exactement quatre exponentiations modulaires sur
le candidat connu, sans crible, sans filtre du Théorème~2 ni
pré-sélection Miller–Rabin.

%------------------------------------------------------------------------------
\subsection{Rang dans les records}
%------------------------------------------------------------------------------

Le Tableau~\ref{tab:tousresultats} (Section~\ref{sec:resultats}) liste
l'ensemble des quatre certificats \textsc{PrimeQuest}. L'instance de
$29\,998$ chiffres est actuellement le plus grand certificat
Pocklington–Lehmer inconditionnel dans la famille
$p = 3m(m+1)+1$, $m = 2^a\cdot 3^b - 1$,
obtenu sur un ordinateur portable grand public sans calcul distribué ni
matériel spécialisé.

%==============================================================================
% PRÉAMBULE MINIMAL AUTONOME (décommenter pour compilation stand-alone)
%==============================================================================
% \iffalse   % <-- supprimer cette ligne pour activer le mode autonome
% \documentclass[11pt,a4paper]{article}
% \usepackage[utf8]{inputenc}
% \usepackage[T1]{fontenc}
% \usepackage[french]{babel}
% \usepackage{amsmath,amssymb,amsthm,mathtools,booktabs}
% \usepackage{geometry}\geometry{margin=2.5cm}
% \usepackage[dvipsnames]{xcolor}
% \usepackage{tcolorbox}
% \tcbuselibrary{skins,breakable}
% \newtcolorbox{rigorousbox}[1][]{colback=Green!8,colframe=Green!60!black,
%   arc=1mm,title=\textbf{Rigoureux},fonttitle=\bfseries,breakable,#1}
% \newtcolorbox{resultbox}[1][]{colback=RoyalBlue!5,colframe=RoyalBlue!70!black,
%   arc=1mm,title=\textbf{Résultat calculatoire},fonttitle=\bfseries,breakable,#1}
% \newtheorem{theorem}{Théorème}[section]
% \newtheorem{lemma}[theorem]{Lemme}
% \theoremstyle{definition}
% \newtheorem{definition}[theorem]{Définition}
% \begin{document}
% %==============================================================================
% Certificat de primalité inconditionnelle — premier de 29 998 chiffres
% Famille paramétrique  p = 3m(m+1)+1,  m = 2^a·3^b − 1
% Vérification numérique indépendante — mai 2026
%==============================================================================
%
% Utilisation dans le document principal :
%   %==============================================================================
% Certificat de primalité inconditionnelle — premier de 29 998 chiffres
% Famille paramétrique  p = 3m(m+1)+1,  m = 2^a·3^b − 1
% Vérification numérique indépendante — mai 2026
%==============================================================================
%
% Utilisation dans le document principal :
%   \input{certificat_29998_fr}
%
% Compilation autonome :
%   pdflatex certificat_29998_fr.tex
%
%==============================================================================

%------------------------------------------------------------------------------
\section{Certificat inconditionnel pour le premier de 29\,998 chiffres}
\label{sec:cert29998}
%------------------------------------------------------------------------------

Cette section présente le certificat de primalité Pocklington–Lehmer
complet et vérifié indépendamment pour le plus grand premier
actuellement produit par \textsc{PrimeQuest}.

%------------------------------------------------------------------------------
\subsection{Paramètres et structure}
%------------------------------------------------------------------------------

\begin{definition}[Instance de 29\,998 chiffres]\label{def:instance29998}
On pose
\[
  a = 19\,435,\qquad b = 19\,173,
\]
et on définit
\[
  m \;=\; 2^{19435}\cdot 3^{19173} - 1
    \qquad\bigl(14\,999 \text{ chiffres décimaux}\bigr),
\]
\[
  p \;=\; 3m(m+1)+1
    \qquad\bigl(29\,998 \text{ chiffres décimaux}\bigr).
\]
\end{definition}

Le rapport $b/a = 19\,173/19\,435 \approx 0{,}9865$ est voisin de l'optimum
théorique $\rho^* = \log_{10}2/\log_{10}3 \approx 0{,}6309$ minimisant le
nombre de chiffres à $a$ fixé ; la valeur retenue ici cible exactement
$29\,998$ chiffres et a été localisée par l'algorithme \textsc{PrimeQuest v4}
(voir Section~\ref{sec:algo}).

\medskip

\noindent\textbf{Identification.}\quad
Les premiers et derniers chiffres décimaux de $p$ sont :
\begin{equation}\label{eq:p29998chiffres}
  p \;=\; 1602087359509060348500037851549319056715\cdots
         16930852697403293697.
\end{equation}

%------------------------------------------------------------------------------
\subsection{Factorisation de $p-1$}
%------------------------------------------------------------------------------

\begin{rigorousbox}
\begin{lemma}[Factorisation de $p-1$]\label{lem:factorisation29998}
Avec $a, b, m, p$ comme dans la Définition~\ref{def:instance29998},
\[
  p - 1 \;=\; \underbrace{2^{19435}\cdot 3^{19174}}_{F}
              \;\cdot\; m,
\]
où $F := 2^{19435}\cdot 3^{19174}$ possède $15\,000$ chiffres décimaux et
vérifie $F > \sqrt{p}$.
\end{lemma}
\end{rigorousbox}

\begin{proof}
Puisque $m + 1 = 2^{19435}\cdot 3^{19173}$,
\[
  p - 1 \;=\; 3m(m+1) \;=\; 3m\cdot 2^{19435}\cdot 3^{19173}
           \;=\; 2^{19435}\cdot 3^{19174}\cdot m.
\]
Pour la borne $F > \sqrt{p}$, on note que $p < 3(m+1)^2$, donc
$\sqrt{p} < \sqrt{3}\,(m+1) < 3(m+1) = 2^{19435}\cdot 3^{19174} = F$.

\noindent Numériquement, en comparant les logarithmes en base $2$ :
\[
  2\log_2 F
  \;=\; 2\bigl(19435 + 19174\log_2 3\bigr)
  \;\approx\; 99\,650{,}14,
\]
\[
  \log_2 p \;\approx\; 99\,648{,}56,
\]
ce qui confirme $F^2 > p$ avec une marge de $\approx 1{,}58$ bit.
\end{proof}

%------------------------------------------------------------------------------
\subsection{Le certificat Pocklington–Lehmer}
%------------------------------------------------------------------------------

Le Théorème~\ref{thm:Pocklington} (rappelé Section~\ref{sec:family})
requiert, pour chaque premier $q \mid F$ — ici $q \in \{2, 3\}$ —, un entier
$w_q$ vérifiant
\[
  w_q^{p-1} \equiv 1 \pmod{p}
  \qquad\text{et}\qquad
  \gcd\!\bigl(w_q^{(p-1)/q} - 1,\; p\bigr) = 1.
\]

\begin{resultbox}[title={\textbf{Certificat} — premier de 29\,998 chiffres}]
\[
  p \;=\; 3m(m+1)+1,\qquad m = 2^{19435}\cdot 3^{19173}-1.
\]

\medskip
\noindent\textbf{Partie connue de $p-1$ :}
$\;F = 2^{19435}\cdot 3^{19174},\quad F > \sqrt{p}.\quad$\checkmark

\medskip
\renewcommand{\arraystretch}{1.4}
\begin{tabular}{@{}c l l@{}}
\toprule
Premier $q$ & Témoin $w_q$ & Conditions vérifiées \\
\midrule
$2$ & $w_2 = 5$ &
  $5^{p-1}\equiv 1\pmod{p}$\quad et\quad
  $\gcd(5^{(p-1)/2}-1,\,p)=1$ \\
$3$ & $w_3 = 7$ &
  $7^{p-1}\equiv 1\pmod{p}$\quad et\quad
  $\gcd(7^{(p-1)/3}-1,\,p)=1$ \\
\bottomrule
\end{tabular}

\medskip
\noindent\textbf{Conclusion :} $p$ est premier (preuve déterministe). \checkmark
\end{resultbox}

%------------------------------------------------------------------------------
\subsection{Vérification numérique indépendante}
%------------------------------------------------------------------------------

Les quatre conditions du certificat ont été vérifiées de façon indépendante
à l'aide de \texttt{Python 3.12} et de son arithmétique entière en précision
arbitraire (backend GMP). Le script ne comporte aucune étape probabiliste ;
toutes les opérations sont des arithmétiques modulaires exactes.

\begin{rigorousbox}[title={\textbf{Attestation de vérification} — calcul indépendant}]
\textbf{Environnement :} Python~3.12,
\texttt{sys.set\_int\_max\_str\_digits(40000)} ;
exponentiations modulaires via \texttt{pow(base, exp, mod)}.\\[4pt]

\textbf{Étape 1 — Reconstruction de $p$.}\quad
Calcul de $m = 2^{19435}\cdot 3^{19173}-1$ (14\,999 chiffres) et de
$p = 3m(m+1)+1$ (29\,998 chiffres) à partir des paramètres bruts.
Premiers 40 chiffres de $p$ :
\texttt{1602087359509060348500037851549319056715} ;
20 derniers : \texttt{16930852697403293697}.
Correspond exactement au certificat stocké.\\[4pt]

\textbf{Étape 2 — Borne Pocklington.}
\[
  2\log_2 F = 99\,650{,}1420 \;>\; \log_2 p = 99\,648{,}5570.\quad\checkmark
\]

\textbf{Étape 3 — Témoin $w_2 = 5$ pour $q = 2$.}
\begin{align*}
  5^{p-1} &\equiv 1 \pmod{p},\quad\checkmark \\
  \gcd(5^{(p-1)/2}-1,\,p) &= 1.\quad\checkmark
\end{align*}

\textbf{Étape 4 — Témoin $w_3 = 7$ pour $q = 3$.}
\begin{align*}
  7^{p-1} &\equiv 1 \pmod{p},\quad\checkmark \\
  \gcd(7^{(p-1)/3}-1,\,p) &= 1.\quad\checkmark
\end{align*}

\textbf{Durée totale :} $6\,945$\,s $\approx 1$\,h\,$56$\,min
(cœur unique, Python GMP, sans parallélisme).\\[4pt]

\textbf{Conclusion :} les quatre conditions Pocklington–Lehmer sont satisfaites.
Le nombre $p$ est \emph{premier de façon inconditionnelle}.
\end{rigorousbox}

%------------------------------------------------------------------------------
\subsection{Comparaison avec le calcul original}
%------------------------------------------------------------------------------

\begin{table}[h]
\centering
\caption{Recherche originale vs.\ vérification indépendante du certificat
         pour le premier de 29\,998 chiffres.}
\label{tab:comparaison29998}
\renewcommand{\arraystretch}{1.3}
\begin{tabular}{@{}lll@{}}
\toprule
 & \textbf{PrimeQuest v4 (recherche)} & \textbf{Vérification indépendante} \\
\midrule
Objectif       & Trouver un candidat premier         & Vérifier le certificat \\
Machine        & ARM64, 9 cœurs (Snapdragon X)       & Cœur unique (Python GMP) \\
Durée          & $11\,146$\,s ($3$\,h\,$06$\,min)   & $6\,945$\,s ($1$\,h\,$56$\,min) \\
Tests MR       & $286$ (pré-filtre probabiliste)     & — (aucune étape probabiliste) \\
Résultat       & Certificat produit                  & Certificat confirmé \\
\bottomrule
\end{tabular}
\end{table}

La vérification est $\approx 1{,}6$ fois plus rapide que la recherche
originale car elle exécute exactement quatre exponentiations modulaires sur
le candidat connu, sans crible, sans filtre du Théorème~2 ni
pré-sélection Miller–Rabin.

%------------------------------------------------------------------------------
\subsection{Rang dans les records}
%------------------------------------------------------------------------------

Le Tableau~\ref{tab:tousresultats} (Section~\ref{sec:resultats}) liste
l'ensemble des quatre certificats \textsc{PrimeQuest}. L'instance de
$29\,998$ chiffres est actuellement le plus grand certificat
Pocklington–Lehmer inconditionnel dans la famille
$p = 3m(m+1)+1$, $m = 2^a\cdot 3^b - 1$,
obtenu sur un ordinateur portable grand public sans calcul distribué ni
matériel spécialisé.

%==============================================================================
% PRÉAMBULE MINIMAL AUTONOME (décommenter pour compilation stand-alone)
%==============================================================================
% \iffalse   % <-- supprimer cette ligne pour activer le mode autonome
% \documentclass[11pt,a4paper]{article}
% \usepackage[utf8]{inputenc}
% \usepackage[T1]{fontenc}
% \usepackage[french]{babel}
% \usepackage{amsmath,amssymb,amsthm,mathtools,booktabs}
% \usepackage{geometry}\geometry{margin=2.5cm}
% \usepackage[dvipsnames]{xcolor}
% \usepackage{tcolorbox}
% \tcbuselibrary{skins,breakable}
% \newtcolorbox{rigorousbox}[1][]{colback=Green!8,colframe=Green!60!black,
%   arc=1mm,title=\textbf{Rigoureux},fonttitle=\bfseries,breakable,#1}
% \newtcolorbox{resultbox}[1][]{colback=RoyalBlue!5,colframe=RoyalBlue!70!black,
%   arc=1mm,title=\textbf{Résultat calculatoire},fonttitle=\bfseries,breakable,#1}
% \newtheorem{theorem}{Théorème}[section]
% \newtheorem{lemma}[theorem]{Lemme}
% \theoremstyle{definition}
% \newtheorem{definition}[theorem]{Définition}
% \begin{document}
% \input{certificat_29998_fr}
% \end{document}
% \fi

%
% Compilation autonome :
%   pdflatex certificat_29998_fr.tex
%
%==============================================================================

%------------------------------------------------------------------------------
\section{Certificat inconditionnel pour le premier de 29\,998 chiffres}
\label{sec:cert29998}
%------------------------------------------------------------------------------

Cette section présente le certificat de primalité Pocklington–Lehmer
complet et vérifié indépendamment pour le plus grand premier
actuellement produit par \textsc{PrimeQuest}.

%------------------------------------------------------------------------------
\subsection{Paramètres et structure}
%------------------------------------------------------------------------------

\begin{definition}[Instance de 29\,998 chiffres]\label{def:instance29998}
On pose
\[
  a = 19\,435,\qquad b = 19\,173,
\]
et on définit
\[
  m \;=\; 2^{19435}\cdot 3^{19173} - 1
    \qquad\bigl(14\,999 \text{ chiffres décimaux}\bigr),
\]
\[
  p \;=\; 3m(m+1)+1
    \qquad\bigl(29\,998 \text{ chiffres décimaux}\bigr).
\]
\end{definition}

Le rapport $b/a = 19\,173/19\,435 \approx 0{,}9865$ est voisin de l'optimum
théorique $\rho^* = \log_{10}2/\log_{10}3 \approx 0{,}6309$ minimisant le
nombre de chiffres à $a$ fixé ; la valeur retenue ici cible exactement
$29\,998$ chiffres et a été localisée par l'algorithme \textsc{PrimeQuest v4}
(voir Section~\ref{sec:algo}).

\medskip

\noindent\textbf{Identification.}\quad
Les premiers et derniers chiffres décimaux de $p$ sont :
\begin{equation}\label{eq:p29998chiffres}
  p \;=\; 1602087359509060348500037851549319056715\cdots
         16930852697403293697.
\end{equation}

%------------------------------------------------------------------------------
\subsection{Factorisation de $p-1$}
%------------------------------------------------------------------------------

\begin{rigorousbox}
\begin{lemma}[Factorisation de $p-1$]\label{lem:factorisation29998}
Avec $a, b, m, p$ comme dans la Définition~\ref{def:instance29998},
\[
  p - 1 \;=\; \underbrace{2^{19435}\cdot 3^{19174}}_{F}
              \;\cdot\; m,
\]
où $F := 2^{19435}\cdot 3^{19174}$ possède $15\,000$ chiffres décimaux et
vérifie $F > \sqrt{p}$.
\end{lemma}
\end{rigorousbox}

\begin{proof}
Puisque $m + 1 = 2^{19435}\cdot 3^{19173}$,
\[
  p - 1 \;=\; 3m(m+1) \;=\; 3m\cdot 2^{19435}\cdot 3^{19173}
           \;=\; 2^{19435}\cdot 3^{19174}\cdot m.
\]
Pour la borne $F > \sqrt{p}$, on note que $p < 3(m+1)^2$, donc
$\sqrt{p} < \sqrt{3}\,(m+1) < 3(m+1) = 2^{19435}\cdot 3^{19174} = F$.

\noindent Numériquement, en comparant les logarithmes en base $2$ :
\[
  2\log_2 F
  \;=\; 2\bigl(19435 + 19174\log_2 3\bigr)
  \;\approx\; 99\,650{,}14,
\]
\[
  \log_2 p \;\approx\; 99\,648{,}56,
\]
ce qui confirme $F^2 > p$ avec une marge de $\approx 1{,}58$ bit.
\end{proof}

%------------------------------------------------------------------------------
\subsection{Le certificat Pocklington–Lehmer}
%------------------------------------------------------------------------------

Le Théorème~\ref{thm:Pocklington} (rappelé Section~\ref{sec:family})
requiert, pour chaque premier $q \mid F$ — ici $q \in \{2, 3\}$ —, un entier
$w_q$ vérifiant
\[
  w_q^{p-1} \equiv 1 \pmod{p}
  \qquad\text{et}\qquad
  \gcd\!\bigl(w_q^{(p-1)/q} - 1,\; p\bigr) = 1.
\]

\begin{resultbox}[title={\textbf{Certificat} — premier de 29\,998 chiffres}]
\[
  p \;=\; 3m(m+1)+1,\qquad m = 2^{19435}\cdot 3^{19173}-1.
\]

\medskip
\noindent\textbf{Partie connue de $p-1$ :}
$\;F = 2^{19435}\cdot 3^{19174},\quad F > \sqrt{p}.\quad$\checkmark

\medskip
\renewcommand{\arraystretch}{1.4}
\begin{tabular}{@{}c l l@{}}
\toprule
Premier $q$ & Témoin $w_q$ & Conditions vérifiées \\
\midrule
$2$ & $w_2 = 5$ &
  $5^{p-1}\equiv 1\pmod{p}$\quad et\quad
  $\gcd(5^{(p-1)/2}-1,\,p)=1$ \\
$3$ & $w_3 = 7$ &
  $7^{p-1}\equiv 1\pmod{p}$\quad et\quad
  $\gcd(7^{(p-1)/3}-1,\,p)=1$ \\
\bottomrule
\end{tabular}

\medskip
\noindent\textbf{Conclusion :} $p$ est premier (preuve déterministe). \checkmark
\end{resultbox}

%------------------------------------------------------------------------------
\subsection{Vérification numérique indépendante}
%------------------------------------------------------------------------------

Les quatre conditions du certificat ont été vérifiées de façon indépendante
à l'aide de \texttt{Python 3.12} et de son arithmétique entière en précision
arbitraire (backend GMP). Le script ne comporte aucune étape probabiliste ;
toutes les opérations sont des arithmétiques modulaires exactes.

\begin{rigorousbox}[title={\textbf{Attestation de vérification} — calcul indépendant}]
\textbf{Environnement :} Python~3.12,
\texttt{sys.set\_int\_max\_str\_digits(40000)} ;
exponentiations modulaires via \texttt{pow(base, exp, mod)}.\\[4pt]

\textbf{Étape 1 — Reconstruction de $p$.}\quad
Calcul de $m = 2^{19435}\cdot 3^{19173}-1$ (14\,999 chiffres) et de
$p = 3m(m+1)+1$ (29\,998 chiffres) à partir des paramètres bruts.
Premiers 40 chiffres de $p$ :
\texttt{1602087359509060348500037851549319056715} ;
20 derniers : \texttt{16930852697403293697}.
Correspond exactement au certificat stocké.\\[4pt]

\textbf{Étape 2 — Borne Pocklington.}
\[
  2\log_2 F = 99\,650{,}1420 \;>\; \log_2 p = 99\,648{,}5570.\quad\checkmark
\]

\textbf{Étape 3 — Témoin $w_2 = 5$ pour $q = 2$.}
\begin{align*}
  5^{p-1} &\equiv 1 \pmod{p},\quad\checkmark \\
  \gcd(5^{(p-1)/2}-1,\,p) &= 1.\quad\checkmark
\end{align*}

\textbf{Étape 4 — Témoin $w_3 = 7$ pour $q = 3$.}
\begin{align*}
  7^{p-1} &\equiv 1 \pmod{p},\quad\checkmark \\
  \gcd(7^{(p-1)/3}-1,\,p) &= 1.\quad\checkmark
\end{align*}

\textbf{Durée totale :} $6\,945$\,s $\approx 1$\,h\,$56$\,min
(cœur unique, Python GMP, sans parallélisme).\\[4pt]

\textbf{Conclusion :} les quatre conditions Pocklington–Lehmer sont satisfaites.
Le nombre $p$ est \emph{premier de façon inconditionnelle}.
\end{rigorousbox}

%------------------------------------------------------------------------------
\subsection{Comparaison avec le calcul original}
%------------------------------------------------------------------------------

\begin{table}[h]
\centering
\caption{Recherche originale vs.\ vérification indépendante du certificat
         pour le premier de 29\,998 chiffres.}
\label{tab:comparaison29998}
\renewcommand{\arraystretch}{1.3}
\begin{tabular}{@{}lll@{}}
\toprule
 & \textbf{PrimeQuest v4 (recherche)} & \textbf{Vérification indépendante} \\
\midrule
Objectif       & Trouver un candidat premier         & Vérifier le certificat \\
Machine        & ARM64, 9 cœurs (Snapdragon X)       & Cœur unique (Python GMP) \\
Durée          & $11\,146$\,s ($3$\,h\,$06$\,min)   & $6\,945$\,s ($1$\,h\,$56$\,min) \\
Tests MR       & $286$ (pré-filtre probabiliste)     & — (aucune étape probabiliste) \\
Résultat       & Certificat produit                  & Certificat confirmé \\
\bottomrule
\end{tabular}
\end{table}

La vérification est $\approx 1{,}6$ fois plus rapide que la recherche
originale car elle exécute exactement quatre exponentiations modulaires sur
le candidat connu, sans crible, sans filtre du Théorème~2 ni
pré-sélection Miller–Rabin.

%------------------------------------------------------------------------------
\subsection{Rang dans les records}
%------------------------------------------------------------------------------

Le Tableau~\ref{tab:tousresultats} (Section~\ref{sec:resultats}) liste
l'ensemble des quatre certificats \textsc{PrimeQuest}. L'instance de
$29\,998$ chiffres est actuellement le plus grand certificat
Pocklington–Lehmer inconditionnel dans la famille
$p = 3m(m+1)+1$, $m = 2^a\cdot 3^b - 1$,
obtenu sur un ordinateur portable grand public sans calcul distribué ni
matériel spécialisé.

%==============================================================================
% PRÉAMBULE MINIMAL AUTONOME (décommenter pour compilation stand-alone)
%==============================================================================
% \iffalse   % <-- supprimer cette ligne pour activer le mode autonome
% \documentclass[11pt,a4paper]{article}
% \usepackage[utf8]{inputenc}
% \usepackage[T1]{fontenc}
% \usepackage[french]{babel}
% \usepackage{amsmath,amssymb,amsthm,mathtools,booktabs}
% \usepackage{geometry}\geometry{margin=2.5cm}
% \usepackage[dvipsnames]{xcolor}
% \usepackage{tcolorbox}
% \tcbuselibrary{skins,breakable}
% \newtcolorbox{rigorousbox}[1][]{colback=Green!8,colframe=Green!60!black,
%   arc=1mm,title=\textbf{Rigoureux},fonttitle=\bfseries,breakable,#1}
% \newtcolorbox{resultbox}[1][]{colback=RoyalBlue!5,colframe=RoyalBlue!70!black,
%   arc=1mm,title=\textbf{Résultat calculatoire},fonttitle=\bfseries,breakable,#1}
% \newtheorem{theorem}{Théorème}[section]
% \newtheorem{lemma}[theorem]{Lemme}
% \theoremstyle{definition}
% \newtheorem{definition}[theorem]{Définition}
% \begin{document}
% %==============================================================================
% Certificat de primalité inconditionnelle — premier de 29 998 chiffres
% Famille paramétrique  p = 3m(m+1)+1,  m = 2^a·3^b − 1
% Vérification numérique indépendante — mai 2026
%==============================================================================
%
% Utilisation dans le document principal :
%   \input{certificat_29998_fr}
%
% Compilation autonome :
%   pdflatex certificat_29998_fr.tex
%
%==============================================================================

%------------------------------------------------------------------------------
\section{Certificat inconditionnel pour le premier de 29\,998 chiffres}
\label{sec:cert29998}
%------------------------------------------------------------------------------

Cette section présente le certificat de primalité Pocklington–Lehmer
complet et vérifié indépendamment pour le plus grand premier
actuellement produit par \textsc{PrimeQuest}.

%------------------------------------------------------------------------------
\subsection{Paramètres et structure}
%------------------------------------------------------------------------------

\begin{definition}[Instance de 29\,998 chiffres]\label{def:instance29998}
On pose
\[
  a = 19\,435,\qquad b = 19\,173,
\]
et on définit
\[
  m \;=\; 2^{19435}\cdot 3^{19173} - 1
    \qquad\bigl(14\,999 \text{ chiffres décimaux}\bigr),
\]
\[
  p \;=\; 3m(m+1)+1
    \qquad\bigl(29\,998 \text{ chiffres décimaux}\bigr).
\]
\end{definition}

Le rapport $b/a = 19\,173/19\,435 \approx 0{,}9865$ est voisin de l'optimum
théorique $\rho^* = \log_{10}2/\log_{10}3 \approx 0{,}6309$ minimisant le
nombre de chiffres à $a$ fixé ; la valeur retenue ici cible exactement
$29\,998$ chiffres et a été localisée par l'algorithme \textsc{PrimeQuest v4}
(voir Section~\ref{sec:algo}).

\medskip

\noindent\textbf{Identification.}\quad
Les premiers et derniers chiffres décimaux de $p$ sont :
\begin{equation}\label{eq:p29998chiffres}
  p \;=\; 1602087359509060348500037851549319056715\cdots
         16930852697403293697.
\end{equation}

%------------------------------------------------------------------------------
\subsection{Factorisation de $p-1$}
%------------------------------------------------------------------------------

\begin{rigorousbox}
\begin{lemma}[Factorisation de $p-1$]\label{lem:factorisation29998}
Avec $a, b, m, p$ comme dans la Définition~\ref{def:instance29998},
\[
  p - 1 \;=\; \underbrace{2^{19435}\cdot 3^{19174}}_{F}
              \;\cdot\; m,
\]
où $F := 2^{19435}\cdot 3^{19174}$ possède $15\,000$ chiffres décimaux et
vérifie $F > \sqrt{p}$.
\end{lemma}
\end{rigorousbox}

\begin{proof}
Puisque $m + 1 = 2^{19435}\cdot 3^{19173}$,
\[
  p - 1 \;=\; 3m(m+1) \;=\; 3m\cdot 2^{19435}\cdot 3^{19173}
           \;=\; 2^{19435}\cdot 3^{19174}\cdot m.
\]
Pour la borne $F > \sqrt{p}$, on note que $p < 3(m+1)^2$, donc
$\sqrt{p} < \sqrt{3}\,(m+1) < 3(m+1) = 2^{19435}\cdot 3^{19174} = F$.

\noindent Numériquement, en comparant les logarithmes en base $2$ :
\[
  2\log_2 F
  \;=\; 2\bigl(19435 + 19174\log_2 3\bigr)
  \;\approx\; 99\,650{,}14,
\]
\[
  \log_2 p \;\approx\; 99\,648{,}56,
\]
ce qui confirme $F^2 > p$ avec une marge de $\approx 1{,}58$ bit.
\end{proof}

%------------------------------------------------------------------------------
\subsection{Le certificat Pocklington–Lehmer}
%------------------------------------------------------------------------------

Le Théorème~\ref{thm:Pocklington} (rappelé Section~\ref{sec:family})
requiert, pour chaque premier $q \mid F$ — ici $q \in \{2, 3\}$ —, un entier
$w_q$ vérifiant
\[
  w_q^{p-1} \equiv 1 \pmod{p}
  \qquad\text{et}\qquad
  \gcd\!\bigl(w_q^{(p-1)/q} - 1,\; p\bigr) = 1.
\]

\begin{resultbox}[title={\textbf{Certificat} — premier de 29\,998 chiffres}]
\[
  p \;=\; 3m(m+1)+1,\qquad m = 2^{19435}\cdot 3^{19173}-1.
\]

\medskip
\noindent\textbf{Partie connue de $p-1$ :}
$\;F = 2^{19435}\cdot 3^{19174},\quad F > \sqrt{p}.\quad$\checkmark

\medskip
\renewcommand{\arraystretch}{1.4}
\begin{tabular}{@{}c l l@{}}
\toprule
Premier $q$ & Témoin $w_q$ & Conditions vérifiées \\
\midrule
$2$ & $w_2 = 5$ &
  $5^{p-1}\equiv 1\pmod{p}$\quad et\quad
  $\gcd(5^{(p-1)/2}-1,\,p)=1$ \\
$3$ & $w_3 = 7$ &
  $7^{p-1}\equiv 1\pmod{p}$\quad et\quad
  $\gcd(7^{(p-1)/3}-1,\,p)=1$ \\
\bottomrule
\end{tabular}

\medskip
\noindent\textbf{Conclusion :} $p$ est premier (preuve déterministe). \checkmark
\end{resultbox}

%------------------------------------------------------------------------------
\subsection{Vérification numérique indépendante}
%------------------------------------------------------------------------------

Les quatre conditions du certificat ont été vérifiées de façon indépendante
à l'aide de \texttt{Python 3.12} et de son arithmétique entière en précision
arbitraire (backend GMP). Le script ne comporte aucune étape probabiliste ;
toutes les opérations sont des arithmétiques modulaires exactes.

\begin{rigorousbox}[title={\textbf{Attestation de vérification} — calcul indépendant}]
\textbf{Environnement :} Python~3.12,
\texttt{sys.set\_int\_max\_str\_digits(40000)} ;
exponentiations modulaires via \texttt{pow(base, exp, mod)}.\\[4pt]

\textbf{Étape 1 — Reconstruction de $p$.}\quad
Calcul de $m = 2^{19435}\cdot 3^{19173}-1$ (14\,999 chiffres) et de
$p = 3m(m+1)+1$ (29\,998 chiffres) à partir des paramètres bruts.
Premiers 40 chiffres de $p$ :
\texttt{1602087359509060348500037851549319056715} ;
20 derniers : \texttt{16930852697403293697}.
Correspond exactement au certificat stocké.\\[4pt]

\textbf{Étape 2 — Borne Pocklington.}
\[
  2\log_2 F = 99\,650{,}1420 \;>\; \log_2 p = 99\,648{,}5570.\quad\checkmark
\]

\textbf{Étape 3 — Témoin $w_2 = 5$ pour $q = 2$.}
\begin{align*}
  5^{p-1} &\equiv 1 \pmod{p},\quad\checkmark \\
  \gcd(5^{(p-1)/2}-1,\,p) &= 1.\quad\checkmark
\end{align*}

\textbf{Étape 4 — Témoin $w_3 = 7$ pour $q = 3$.}
\begin{align*}
  7^{p-1} &\equiv 1 \pmod{p},\quad\checkmark \\
  \gcd(7^{(p-1)/3}-1,\,p) &= 1.\quad\checkmark
\end{align*}

\textbf{Durée totale :} $6\,945$\,s $\approx 1$\,h\,$56$\,min
(cœur unique, Python GMP, sans parallélisme).\\[4pt]

\textbf{Conclusion :} les quatre conditions Pocklington–Lehmer sont satisfaites.
Le nombre $p$ est \emph{premier de façon inconditionnelle}.
\end{rigorousbox}

%------------------------------------------------------------------------------
\subsection{Comparaison avec le calcul original}
%------------------------------------------------------------------------------

\begin{table}[h]
\centering
\caption{Recherche originale vs.\ vérification indépendante du certificat
         pour le premier de 29\,998 chiffres.}
\label{tab:comparaison29998}
\renewcommand{\arraystretch}{1.3}
\begin{tabular}{@{}lll@{}}
\toprule
 & \textbf{PrimeQuest v4 (recherche)} & \textbf{Vérification indépendante} \\
\midrule
Objectif       & Trouver un candidat premier         & Vérifier le certificat \\
Machine        & ARM64, 9 cœurs (Snapdragon X)       & Cœur unique (Python GMP) \\
Durée          & $11\,146$\,s ($3$\,h\,$06$\,min)   & $6\,945$\,s ($1$\,h\,$56$\,min) \\
Tests MR       & $286$ (pré-filtre probabiliste)     & — (aucune étape probabiliste) \\
Résultat       & Certificat produit                  & Certificat confirmé \\
\bottomrule
\end{tabular}
\end{table}

La vérification est $\approx 1{,}6$ fois plus rapide que la recherche
originale car elle exécute exactement quatre exponentiations modulaires sur
le candidat connu, sans crible, sans filtre du Théorème~2 ni
pré-sélection Miller–Rabin.

%------------------------------------------------------------------------------
\subsection{Rang dans les records}
%------------------------------------------------------------------------------

Le Tableau~\ref{tab:tousresultats} (Section~\ref{sec:resultats}) liste
l'ensemble des quatre certificats \textsc{PrimeQuest}. L'instance de
$29\,998$ chiffres est actuellement le plus grand certificat
Pocklington–Lehmer inconditionnel dans la famille
$p = 3m(m+1)+1$, $m = 2^a\cdot 3^b - 1$,
obtenu sur un ordinateur portable grand public sans calcul distribué ni
matériel spécialisé.

%==============================================================================
% PRÉAMBULE MINIMAL AUTONOME (décommenter pour compilation stand-alone)
%==============================================================================
% \iffalse   % <-- supprimer cette ligne pour activer le mode autonome
% \documentclass[11pt,a4paper]{article}
% \usepackage[utf8]{inputenc}
% \usepackage[T1]{fontenc}
% \usepackage[french]{babel}
% \usepackage{amsmath,amssymb,amsthm,mathtools,booktabs}
% \usepackage{geometry}\geometry{margin=2.5cm}
% \usepackage[dvipsnames]{xcolor}
% \usepackage{tcolorbox}
% \tcbuselibrary{skins,breakable}
% \newtcolorbox{rigorousbox}[1][]{colback=Green!8,colframe=Green!60!black,
%   arc=1mm,title=\textbf{Rigoureux},fonttitle=\bfseries,breakable,#1}
% \newtcolorbox{resultbox}[1][]{colback=RoyalBlue!5,colframe=RoyalBlue!70!black,
%   arc=1mm,title=\textbf{Résultat calculatoire},fonttitle=\bfseries,breakable,#1}
% \newtheorem{theorem}{Théorème}[section]
% \newtheorem{lemma}[theorem]{Lemme}
% \theoremstyle{definition}
% \newtheorem{definition}[theorem]{Définition}
% \begin{document}
% \input{certificat_29998_fr}
% \end{document}
% \fi

% \end{document}
% \fi

% \end{document}
% \fi

%
% Compilation autonome :
%   pdflatex certificat_29998_fr.tex
%
%==============================================================================

%------------------------------------------------------------------------------
\section{Certificat inconditionnel pour le premier de 29\,998 chiffres}
\label{sec:cert29998}
%------------------------------------------------------------------------------

Cette section présente le certificat de primalité Pocklington–Lehmer
complet et vérifié indépendamment pour le plus grand premier
actuellement produit par \textsc{PrimeQuest}.

%------------------------------------------------------------------------------
\subsection{Paramètres et structure}
%------------------------------------------------------------------------------

\begin{definition}[Instance de 29\,998 chiffres]\label{def:instance29998}
On pose
\[
  a = 19\,435,\qquad b = 19\,173,
\]
et on définit
\[
  m \;=\; 2^{19435}\cdot 3^{19173} - 1
    \qquad\bigl(14\,999 \text{ chiffres décimaux}\bigr),
\]
\[
  p \;=\; 3m(m+1)+1
    \qquad\bigl(29\,998 \text{ chiffres décimaux}\bigr).
\]
\end{definition}

Le rapport $b/a = 19\,173/19\,435 \approx 0{,}9865$ est voisin de l'optimum
théorique $\rho^* = \log_{10}2/\log_{10}3 \approx 0{,}6309$ minimisant le
nombre de chiffres à $a$ fixé ; la valeur retenue ici cible exactement
$29\,998$ chiffres et a été localisée par l'algorithme \textsc{PrimeQuest v4}
(voir Section~\ref{sec:algo}).

\medskip

\noindent\textbf{Identification.}\quad
Les premiers et derniers chiffres décimaux de $p$ sont :
\begin{equation}\label{eq:p29998chiffres}
  p \;=\; 1602087359509060348500037851549319056715\cdots
         16930852697403293697.
\end{equation}

%------------------------------------------------------------------------------
\subsection{Factorisation de $p-1$}
%------------------------------------------------------------------------------

\begin{rigorousbox}
\begin{lemma}[Factorisation de $p-1$]\label{lem:factorisation29998}
Avec $a, b, m, p$ comme dans la Définition~\ref{def:instance29998},
\[
  p - 1 \;=\; \underbrace{2^{19435}\cdot 3^{19174}}_{F}
              \;\cdot\; m,
\]
où $F := 2^{19435}\cdot 3^{19174}$ possède $15\,000$ chiffres décimaux et
vérifie $F > \sqrt{p}$.
\end{lemma}
\end{rigorousbox}

\begin{proof}
Puisque $m + 1 = 2^{19435}\cdot 3^{19173}$,
\[
  p - 1 \;=\; 3m(m+1) \;=\; 3m\cdot 2^{19435}\cdot 3^{19173}
           \;=\; 2^{19435}\cdot 3^{19174}\cdot m.
\]
Pour la borne $F > \sqrt{p}$, on note que $p < 3(m+1)^2$, donc
$\sqrt{p} < \sqrt{3}\,(m+1) < 3(m+1) = 2^{19435}\cdot 3^{19174} = F$.

\noindent Numériquement, en comparant les logarithmes en base $2$ :
\[
  2\log_2 F
  \;=\; 2\bigl(19435 + 19174\log_2 3\bigr)
  \;\approx\; 99\,650{,}14,
\]
\[
  \log_2 p \;\approx\; 99\,648{,}56,
\]
ce qui confirme $F^2 > p$ avec une marge de $\approx 1{,}58$ bit.
\end{proof}

%------------------------------------------------------------------------------
\subsection{Le certificat Pocklington–Lehmer}
%------------------------------------------------------------------------------

Le Théorème~\ref{thm:Pocklington} (rappelé Section~\ref{sec:family})
requiert, pour chaque premier $q \mid F$ — ici $q \in \{2, 3\}$ —, un entier
$w_q$ vérifiant
\[
  w_q^{p-1} \equiv 1 \pmod{p}
  \qquad\text{et}\qquad
  \gcd\!\bigl(w_q^{(p-1)/q} - 1,\; p\bigr) = 1.
\]

\begin{resultbox}[title={\textbf{Certificat} — premier de 29\,998 chiffres}]
\[
  p \;=\; 3m(m+1)+1,\qquad m = 2^{19435}\cdot 3^{19173}-1.
\]

\medskip
\noindent\textbf{Partie connue de $p-1$ :}
$\;F = 2^{19435}\cdot 3^{19174},\quad F > \sqrt{p}.\quad$\checkmark

\medskip
\renewcommand{\arraystretch}{1.4}
\begin{tabular}{@{}c l l@{}}
\toprule
Premier $q$ & Témoin $w_q$ & Conditions vérifiées \\
\midrule
$2$ & $w_2 = 5$ &
  $5^{p-1}\equiv 1\pmod{p}$\quad et\quad
  $\gcd(5^{(p-1)/2}-1,\,p)=1$ \\
$3$ & $w_3 = 7$ &
  $7^{p-1}\equiv 1\pmod{p}$\quad et\quad
  $\gcd(7^{(p-1)/3}-1,\,p)=1$ \\
\bottomrule
\end{tabular}

\medskip
\noindent\textbf{Conclusion :} $p$ est premier (preuve déterministe). \checkmark
\end{resultbox}

%------------------------------------------------------------------------------
\subsection{Vérification numérique indépendante}
%------------------------------------------------------------------------------

Les quatre conditions du certificat ont été vérifiées de façon indépendante
à l'aide de \texttt{Python 3.12} et de son arithmétique entière en précision
arbitraire (backend GMP). Le script ne comporte aucune étape probabiliste ;
toutes les opérations sont des arithmétiques modulaires exactes.

\begin{rigorousbox}[title={\textbf{Attestation de vérification} — calcul indépendant}]
\textbf{Environnement :} Python~3.12,
\texttt{sys.set\_int\_max\_str\_digits(40000)} ;
exponentiations modulaires via \texttt{pow(base, exp, mod)}.\\[4pt]

\textbf{Étape 1 — Reconstruction de $p$.}\quad
Calcul de $m = 2^{19435}\cdot 3^{19173}-1$ (14\,999 chiffres) et de
$p = 3m(m+1)+1$ (29\,998 chiffres) à partir des paramètres bruts.
Premiers 40 chiffres de $p$ :
\texttt{1602087359509060348500037851549319056715} ;
20 derniers : \texttt{16930852697403293697}.
Correspond exactement au certificat stocké.\\[4pt]

\textbf{Étape 2 — Borne Pocklington.}
\[
  2\log_2 F = 99\,650{,}1420 \;>\; \log_2 p = 99\,648{,}5570.\quad\checkmark
\]

\textbf{Étape 3 — Témoin $w_2 = 5$ pour $q = 2$.}
\begin{align*}
  5^{p-1} &\equiv 1 \pmod{p},\quad\checkmark \\
  \gcd(5^{(p-1)/2}-1,\,p) &= 1.\quad\checkmark
\end{align*}

\textbf{Étape 4 — Témoin $w_3 = 7$ pour $q = 3$.}
\begin{align*}
  7^{p-1} &\equiv 1 \pmod{p},\quad\checkmark \\
  \gcd(7^{(p-1)/3}-1,\,p) &= 1.\quad\checkmark
\end{align*}

\textbf{Durée totale :} $6\,945$\,s $\approx 1$\,h\,$56$\,min
(cœur unique, Python GMP, sans parallélisme).\\[4pt]

\textbf{Conclusion :} les quatre conditions Pocklington–Lehmer sont satisfaites.
Le nombre $p$ est \emph{premier de façon inconditionnelle}.
\end{rigorousbox}

%------------------------------------------------------------------------------
\subsection{Comparaison avec le calcul original}
%------------------------------------------------------------------------------

\begin{table}[h]
\centering
\caption{Recherche originale vs.\ vérification indépendante du certificat
         pour le premier de 29\,998 chiffres.}
\label{tab:comparaison29998}
\renewcommand{\arraystretch}{1.3}
\begin{tabular}{@{}lll@{}}
\toprule
 & \textbf{PrimeQuest v4 (recherche)} & \textbf{Vérification indépendante} \\
\midrule
Objectif       & Trouver un candidat premier         & Vérifier le certificat \\
Machine        & ARM64, 9 cœurs (Snapdragon X)       & Cœur unique (Python GMP) \\
Durée          & $11\,146$\,s ($3$\,h\,$06$\,min)   & $6\,945$\,s ($1$\,h\,$56$\,min) \\
Tests MR       & $286$ (pré-filtre probabiliste)     & — (aucune étape probabiliste) \\
Résultat       & Certificat produit                  & Certificat confirmé \\
\bottomrule
\end{tabular}
\end{table}

La vérification est $\approx 1{,}6$ fois plus rapide que la recherche
originale car elle exécute exactement quatre exponentiations modulaires sur
le candidat connu, sans crible, sans filtre du Théorème~2 ni
pré-sélection Miller–Rabin.

%------------------------------------------------------------------------------
\subsection{Rang dans les records}
%------------------------------------------------------------------------------

Le Tableau~\ref{tab:tousresultats} (Section~\ref{sec:resultats}) liste
l'ensemble des quatre certificats \textsc{PrimeQuest}. L'instance de
$29\,998$ chiffres est actuellement le plus grand certificat
Pocklington–Lehmer inconditionnel dans la famille
$p = 3m(m+1)+1$, $m = 2^a\cdot 3^b - 1$,
obtenu sur un ordinateur portable grand public sans calcul distribué ni
matériel spécialisé.

%==============================================================================
% PRÉAMBULE MINIMAL AUTONOME (décommenter pour compilation stand-alone)
%==============================================================================
% \iffalse   % <-- supprimer cette ligne pour activer le mode autonome
% \documentclass[11pt,a4paper]{article}
% \usepackage[utf8]{inputenc}
% \usepackage[T1]{fontenc}
% \usepackage[french]{babel}
% \usepackage{amsmath,amssymb,amsthm,mathtools,booktabs}
% \usepackage{geometry}\geometry{margin=2.5cm}
% \usepackage[dvipsnames]{xcolor}
% \usepackage{tcolorbox}
% \tcbuselibrary{skins,breakable}
% \newtcolorbox{rigorousbox}[1][]{colback=Green!8,colframe=Green!60!black,
%   arc=1mm,title=\textbf{Rigoureux},fonttitle=\bfseries,breakable,#1}
% \newtcolorbox{resultbox}[1][]{colback=RoyalBlue!5,colframe=RoyalBlue!70!black,
%   arc=1mm,title=\textbf{Résultat calculatoire},fonttitle=\bfseries,breakable,#1}
% \newtheorem{theorem}{Théorème}[section]
% \newtheorem{lemma}[theorem]{Lemme}
% \theoremstyle{definition}
% \newtheorem{definition}[theorem]{Définition}
% \begin{document}
% %==============================================================================
% Certificat de primalité inconditionnelle — premier de 29 998 chiffres
% Famille paramétrique  p = 3m(m+1)+1,  m = 2^a·3^b − 1
% Vérification numérique indépendante — mai 2026
%==============================================================================
%
% Utilisation dans le document principal :
%   %==============================================================================
% Certificat de primalité inconditionnelle — premier de 29 998 chiffres
% Famille paramétrique  p = 3m(m+1)+1,  m = 2^a·3^b − 1
% Vérification numérique indépendante — mai 2026
%==============================================================================
%
% Utilisation dans le document principal :
%   %==============================================================================
% Certificat de primalité inconditionnelle — premier de 29 998 chiffres
% Famille paramétrique  p = 3m(m+1)+1,  m = 2^a·3^b − 1
% Vérification numérique indépendante — mai 2026
%==============================================================================
%
% Utilisation dans le document principal :
%   \input{certificat_29998_fr}
%
% Compilation autonome :
%   pdflatex certificat_29998_fr.tex
%
%==============================================================================

%------------------------------------------------------------------------------
\section{Certificat inconditionnel pour le premier de 29\,998 chiffres}
\label{sec:cert29998}
%------------------------------------------------------------------------------

Cette section présente le certificat de primalité Pocklington–Lehmer
complet et vérifié indépendamment pour le plus grand premier
actuellement produit par \textsc{PrimeQuest}.

%------------------------------------------------------------------------------
\subsection{Paramètres et structure}
%------------------------------------------------------------------------------

\begin{definition}[Instance de 29\,998 chiffres]\label{def:instance29998}
On pose
\[
  a = 19\,435,\qquad b = 19\,173,
\]
et on définit
\[
  m \;=\; 2^{19435}\cdot 3^{19173} - 1
    \qquad\bigl(14\,999 \text{ chiffres décimaux}\bigr),
\]
\[
  p \;=\; 3m(m+1)+1
    \qquad\bigl(29\,998 \text{ chiffres décimaux}\bigr).
\]
\end{definition}

Le rapport $b/a = 19\,173/19\,435 \approx 0{,}9865$ est voisin de l'optimum
théorique $\rho^* = \log_{10}2/\log_{10}3 \approx 0{,}6309$ minimisant le
nombre de chiffres à $a$ fixé ; la valeur retenue ici cible exactement
$29\,998$ chiffres et a été localisée par l'algorithme \textsc{PrimeQuest v4}
(voir Section~\ref{sec:algo}).

\medskip

\noindent\textbf{Identification.}\quad
Les premiers et derniers chiffres décimaux de $p$ sont :
\begin{equation}\label{eq:p29998chiffres}
  p \;=\; 1602087359509060348500037851549319056715\cdots
         16930852697403293697.
\end{equation}

%------------------------------------------------------------------------------
\subsection{Factorisation de $p-1$}
%------------------------------------------------------------------------------

\begin{rigorousbox}
\begin{lemma}[Factorisation de $p-1$]\label{lem:factorisation29998}
Avec $a, b, m, p$ comme dans la Définition~\ref{def:instance29998},
\[
  p - 1 \;=\; \underbrace{2^{19435}\cdot 3^{19174}}_{F}
              \;\cdot\; m,
\]
où $F := 2^{19435}\cdot 3^{19174}$ possède $15\,000$ chiffres décimaux et
vérifie $F > \sqrt{p}$.
\end{lemma}
\end{rigorousbox}

\begin{proof}
Puisque $m + 1 = 2^{19435}\cdot 3^{19173}$,
\[
  p - 1 \;=\; 3m(m+1) \;=\; 3m\cdot 2^{19435}\cdot 3^{19173}
           \;=\; 2^{19435}\cdot 3^{19174}\cdot m.
\]
Pour la borne $F > \sqrt{p}$, on note que $p < 3(m+1)^2$, donc
$\sqrt{p} < \sqrt{3}\,(m+1) < 3(m+1) = 2^{19435}\cdot 3^{19174} = F$.

\noindent Numériquement, en comparant les logarithmes en base $2$ :
\[
  2\log_2 F
  \;=\; 2\bigl(19435 + 19174\log_2 3\bigr)
  \;\approx\; 99\,650{,}14,
\]
\[
  \log_2 p \;\approx\; 99\,648{,}56,
\]
ce qui confirme $F^2 > p$ avec une marge de $\approx 1{,}58$ bit.
\end{proof}

%------------------------------------------------------------------------------
\subsection{Le certificat Pocklington–Lehmer}
%------------------------------------------------------------------------------

Le Théorème~\ref{thm:Pocklington} (rappelé Section~\ref{sec:family})
requiert, pour chaque premier $q \mid F$ — ici $q \in \{2, 3\}$ —, un entier
$w_q$ vérifiant
\[
  w_q^{p-1} \equiv 1 \pmod{p}
  \qquad\text{et}\qquad
  \gcd\!\bigl(w_q^{(p-1)/q} - 1,\; p\bigr) = 1.
\]

\begin{resultbox}[title={\textbf{Certificat} — premier de 29\,998 chiffres}]
\[
  p \;=\; 3m(m+1)+1,\qquad m = 2^{19435}\cdot 3^{19173}-1.
\]

\medskip
\noindent\textbf{Partie connue de $p-1$ :}
$\;F = 2^{19435}\cdot 3^{19174},\quad F > \sqrt{p}.\quad$\checkmark

\medskip
\renewcommand{\arraystretch}{1.4}
\begin{tabular}{@{}c l l@{}}
\toprule
Premier $q$ & Témoin $w_q$ & Conditions vérifiées \\
\midrule
$2$ & $w_2 = 5$ &
  $5^{p-1}\equiv 1\pmod{p}$\quad et\quad
  $\gcd(5^{(p-1)/2}-1,\,p)=1$ \\
$3$ & $w_3 = 7$ &
  $7^{p-1}\equiv 1\pmod{p}$\quad et\quad
  $\gcd(7^{(p-1)/3}-1,\,p)=1$ \\
\bottomrule
\end{tabular}

\medskip
\noindent\textbf{Conclusion :} $p$ est premier (preuve déterministe). \checkmark
\end{resultbox}

%------------------------------------------------------------------------------
\subsection{Vérification numérique indépendante}
%------------------------------------------------------------------------------

Les quatre conditions du certificat ont été vérifiées de façon indépendante
à l'aide de \texttt{Python 3.12} et de son arithmétique entière en précision
arbitraire (backend GMP). Le script ne comporte aucune étape probabiliste ;
toutes les opérations sont des arithmétiques modulaires exactes.

\begin{rigorousbox}[title={\textbf{Attestation de vérification} — calcul indépendant}]
\textbf{Environnement :} Python~3.12,
\texttt{sys.set\_int\_max\_str\_digits(40000)} ;
exponentiations modulaires via \texttt{pow(base, exp, mod)}.\\[4pt]

\textbf{Étape 1 — Reconstruction de $p$.}\quad
Calcul de $m = 2^{19435}\cdot 3^{19173}-1$ (14\,999 chiffres) et de
$p = 3m(m+1)+1$ (29\,998 chiffres) à partir des paramètres bruts.
Premiers 40 chiffres de $p$ :
\texttt{1602087359509060348500037851549319056715} ;
20 derniers : \texttt{16930852697403293697}.
Correspond exactement au certificat stocké.\\[4pt]

\textbf{Étape 2 — Borne Pocklington.}
\[
  2\log_2 F = 99\,650{,}1420 \;>\; \log_2 p = 99\,648{,}5570.\quad\checkmark
\]

\textbf{Étape 3 — Témoin $w_2 = 5$ pour $q = 2$.}
\begin{align*}
  5^{p-1} &\equiv 1 \pmod{p},\quad\checkmark \\
  \gcd(5^{(p-1)/2}-1,\,p) &= 1.\quad\checkmark
\end{align*}

\textbf{Étape 4 — Témoin $w_3 = 7$ pour $q = 3$.}
\begin{align*}
  7^{p-1} &\equiv 1 \pmod{p},\quad\checkmark \\
  \gcd(7^{(p-1)/3}-1,\,p) &= 1.\quad\checkmark
\end{align*}

\textbf{Durée totale :} $6\,945$\,s $\approx 1$\,h\,$56$\,min
(cœur unique, Python GMP, sans parallélisme).\\[4pt]

\textbf{Conclusion :} les quatre conditions Pocklington–Lehmer sont satisfaites.
Le nombre $p$ est \emph{premier de façon inconditionnelle}.
\end{rigorousbox}

%------------------------------------------------------------------------------
\subsection{Comparaison avec le calcul original}
%------------------------------------------------------------------------------

\begin{table}[h]
\centering
\caption{Recherche originale vs.\ vérification indépendante du certificat
         pour le premier de 29\,998 chiffres.}
\label{tab:comparaison29998}
\renewcommand{\arraystretch}{1.3}
\begin{tabular}{@{}lll@{}}
\toprule
 & \textbf{PrimeQuest v4 (recherche)} & \textbf{Vérification indépendante} \\
\midrule
Objectif       & Trouver un candidat premier         & Vérifier le certificat \\
Machine        & ARM64, 9 cœurs (Snapdragon X)       & Cœur unique (Python GMP) \\
Durée          & $11\,146$\,s ($3$\,h\,$06$\,min)   & $6\,945$\,s ($1$\,h\,$56$\,min) \\
Tests MR       & $286$ (pré-filtre probabiliste)     & — (aucune étape probabiliste) \\
Résultat       & Certificat produit                  & Certificat confirmé \\
\bottomrule
\end{tabular}
\end{table}

La vérification est $\approx 1{,}6$ fois plus rapide que la recherche
originale car elle exécute exactement quatre exponentiations modulaires sur
le candidat connu, sans crible, sans filtre du Théorème~2 ni
pré-sélection Miller–Rabin.

%------------------------------------------------------------------------------
\subsection{Rang dans les records}
%------------------------------------------------------------------------------

Le Tableau~\ref{tab:tousresultats} (Section~\ref{sec:resultats}) liste
l'ensemble des quatre certificats \textsc{PrimeQuest}. L'instance de
$29\,998$ chiffres est actuellement le plus grand certificat
Pocklington–Lehmer inconditionnel dans la famille
$p = 3m(m+1)+1$, $m = 2^a\cdot 3^b - 1$,
obtenu sur un ordinateur portable grand public sans calcul distribué ni
matériel spécialisé.

%==============================================================================
% PRÉAMBULE MINIMAL AUTONOME (décommenter pour compilation stand-alone)
%==============================================================================
% \iffalse   % <-- supprimer cette ligne pour activer le mode autonome
% \documentclass[11pt,a4paper]{article}
% \usepackage[utf8]{inputenc}
% \usepackage[T1]{fontenc}
% \usepackage[french]{babel}
% \usepackage{amsmath,amssymb,amsthm,mathtools,booktabs}
% \usepackage{geometry}\geometry{margin=2.5cm}
% \usepackage[dvipsnames]{xcolor}
% \usepackage{tcolorbox}
% \tcbuselibrary{skins,breakable}
% \newtcolorbox{rigorousbox}[1][]{colback=Green!8,colframe=Green!60!black,
%   arc=1mm,title=\textbf{Rigoureux},fonttitle=\bfseries,breakable,#1}
% \newtcolorbox{resultbox}[1][]{colback=RoyalBlue!5,colframe=RoyalBlue!70!black,
%   arc=1mm,title=\textbf{Résultat calculatoire},fonttitle=\bfseries,breakable,#1}
% \newtheorem{theorem}{Théorème}[section]
% \newtheorem{lemma}[theorem]{Lemme}
% \theoremstyle{definition}
% \newtheorem{definition}[theorem]{Définition}
% \begin{document}
% \input{certificat_29998_fr}
% \end{document}
% \fi

%
% Compilation autonome :
%   pdflatex certificat_29998_fr.tex
%
%==============================================================================

%------------------------------------------------------------------------------
\section{Certificat inconditionnel pour le premier de 29\,998 chiffres}
\label{sec:cert29998}
%------------------------------------------------------------------------------

Cette section présente le certificat de primalité Pocklington–Lehmer
complet et vérifié indépendamment pour le plus grand premier
actuellement produit par \textsc{PrimeQuest}.

%------------------------------------------------------------------------------
\subsection{Paramètres et structure}
%------------------------------------------------------------------------------

\begin{definition}[Instance de 29\,998 chiffres]\label{def:instance29998}
On pose
\[
  a = 19\,435,\qquad b = 19\,173,
\]
et on définit
\[
  m \;=\; 2^{19435}\cdot 3^{19173} - 1
    \qquad\bigl(14\,999 \text{ chiffres décimaux}\bigr),
\]
\[
  p \;=\; 3m(m+1)+1
    \qquad\bigl(29\,998 \text{ chiffres décimaux}\bigr).
\]
\end{definition}

Le rapport $b/a = 19\,173/19\,435 \approx 0{,}9865$ est voisin de l'optimum
théorique $\rho^* = \log_{10}2/\log_{10}3 \approx 0{,}6309$ minimisant le
nombre de chiffres à $a$ fixé ; la valeur retenue ici cible exactement
$29\,998$ chiffres et a été localisée par l'algorithme \textsc{PrimeQuest v4}
(voir Section~\ref{sec:algo}).

\medskip

\noindent\textbf{Identification.}\quad
Les premiers et derniers chiffres décimaux de $p$ sont :
\begin{equation}\label{eq:p29998chiffres}
  p \;=\; 1602087359509060348500037851549319056715\cdots
         16930852697403293697.
\end{equation}

%------------------------------------------------------------------------------
\subsection{Factorisation de $p-1$}
%------------------------------------------------------------------------------

\begin{rigorousbox}
\begin{lemma}[Factorisation de $p-1$]\label{lem:factorisation29998}
Avec $a, b, m, p$ comme dans la Définition~\ref{def:instance29998},
\[
  p - 1 \;=\; \underbrace{2^{19435}\cdot 3^{19174}}_{F}
              \;\cdot\; m,
\]
où $F := 2^{19435}\cdot 3^{19174}$ possède $15\,000$ chiffres décimaux et
vérifie $F > \sqrt{p}$.
\end{lemma}
\end{rigorousbox}

\begin{proof}
Puisque $m + 1 = 2^{19435}\cdot 3^{19173}$,
\[
  p - 1 \;=\; 3m(m+1) \;=\; 3m\cdot 2^{19435}\cdot 3^{19173}
           \;=\; 2^{19435}\cdot 3^{19174}\cdot m.
\]
Pour la borne $F > \sqrt{p}$, on note que $p < 3(m+1)^2$, donc
$\sqrt{p} < \sqrt{3}\,(m+1) < 3(m+1) = 2^{19435}\cdot 3^{19174} = F$.

\noindent Numériquement, en comparant les logarithmes en base $2$ :
\[
  2\log_2 F
  \;=\; 2\bigl(19435 + 19174\log_2 3\bigr)
  \;\approx\; 99\,650{,}14,
\]
\[
  \log_2 p \;\approx\; 99\,648{,}56,
\]
ce qui confirme $F^2 > p$ avec une marge de $\approx 1{,}58$ bit.
\end{proof}

%------------------------------------------------------------------------------
\subsection{Le certificat Pocklington–Lehmer}
%------------------------------------------------------------------------------

Le Théorème~\ref{thm:Pocklington} (rappelé Section~\ref{sec:family})
requiert, pour chaque premier $q \mid F$ — ici $q \in \{2, 3\}$ —, un entier
$w_q$ vérifiant
\[
  w_q^{p-1} \equiv 1 \pmod{p}
  \qquad\text{et}\qquad
  \gcd\!\bigl(w_q^{(p-1)/q} - 1,\; p\bigr) = 1.
\]

\begin{resultbox}[title={\textbf{Certificat} — premier de 29\,998 chiffres}]
\[
  p \;=\; 3m(m+1)+1,\qquad m = 2^{19435}\cdot 3^{19173}-1.
\]

\medskip
\noindent\textbf{Partie connue de $p-1$ :}
$\;F = 2^{19435}\cdot 3^{19174},\quad F > \sqrt{p}.\quad$\checkmark

\medskip
\renewcommand{\arraystretch}{1.4}
\begin{tabular}{@{}c l l@{}}
\toprule
Premier $q$ & Témoin $w_q$ & Conditions vérifiées \\
\midrule
$2$ & $w_2 = 5$ &
  $5^{p-1}\equiv 1\pmod{p}$\quad et\quad
  $\gcd(5^{(p-1)/2}-1,\,p)=1$ \\
$3$ & $w_3 = 7$ &
  $7^{p-1}\equiv 1\pmod{p}$\quad et\quad
  $\gcd(7^{(p-1)/3}-1,\,p)=1$ \\
\bottomrule
\end{tabular}

\medskip
\noindent\textbf{Conclusion :} $p$ est premier (preuve déterministe). \checkmark
\end{resultbox}

%------------------------------------------------------------------------------
\subsection{Vérification numérique indépendante}
%------------------------------------------------------------------------------

Les quatre conditions du certificat ont été vérifiées de façon indépendante
à l'aide de \texttt{Python 3.12} et de son arithmétique entière en précision
arbitraire (backend GMP). Le script ne comporte aucune étape probabiliste ;
toutes les opérations sont des arithmétiques modulaires exactes.

\begin{rigorousbox}[title={\textbf{Attestation de vérification} — calcul indépendant}]
\textbf{Environnement :} Python~3.12,
\texttt{sys.set\_int\_max\_str\_digits(40000)} ;
exponentiations modulaires via \texttt{pow(base, exp, mod)}.\\[4pt]

\textbf{Étape 1 — Reconstruction de $p$.}\quad
Calcul de $m = 2^{19435}\cdot 3^{19173}-1$ (14\,999 chiffres) et de
$p = 3m(m+1)+1$ (29\,998 chiffres) à partir des paramètres bruts.
Premiers 40 chiffres de $p$ :
\texttt{1602087359509060348500037851549319056715} ;
20 derniers : \texttt{16930852697403293697}.
Correspond exactement au certificat stocké.\\[4pt]

\textbf{Étape 2 — Borne Pocklington.}
\[
  2\log_2 F = 99\,650{,}1420 \;>\; \log_2 p = 99\,648{,}5570.\quad\checkmark
\]

\textbf{Étape 3 — Témoin $w_2 = 5$ pour $q = 2$.}
\begin{align*}
  5^{p-1} &\equiv 1 \pmod{p},\quad\checkmark \\
  \gcd(5^{(p-1)/2}-1,\,p) &= 1.\quad\checkmark
\end{align*}

\textbf{Étape 4 — Témoin $w_3 = 7$ pour $q = 3$.}
\begin{align*}
  7^{p-1} &\equiv 1 \pmod{p},\quad\checkmark \\
  \gcd(7^{(p-1)/3}-1,\,p) &= 1.\quad\checkmark
\end{align*}

\textbf{Durée totale :} $6\,945$\,s $\approx 1$\,h\,$56$\,min
(cœur unique, Python GMP, sans parallélisme).\\[4pt]

\textbf{Conclusion :} les quatre conditions Pocklington–Lehmer sont satisfaites.
Le nombre $p$ est \emph{premier de façon inconditionnelle}.
\end{rigorousbox}

%------------------------------------------------------------------------------
\subsection{Comparaison avec le calcul original}
%------------------------------------------------------------------------------

\begin{table}[h]
\centering
\caption{Recherche originale vs.\ vérification indépendante du certificat
         pour le premier de 29\,998 chiffres.}
\label{tab:comparaison29998}
\renewcommand{\arraystretch}{1.3}
\begin{tabular}{@{}lll@{}}
\toprule
 & \textbf{PrimeQuest v4 (recherche)} & \textbf{Vérification indépendante} \\
\midrule
Objectif       & Trouver un candidat premier         & Vérifier le certificat \\
Machine        & ARM64, 9 cœurs (Snapdragon X)       & Cœur unique (Python GMP) \\
Durée          & $11\,146$\,s ($3$\,h\,$06$\,min)   & $6\,945$\,s ($1$\,h\,$56$\,min) \\
Tests MR       & $286$ (pré-filtre probabiliste)     & — (aucune étape probabiliste) \\
Résultat       & Certificat produit                  & Certificat confirmé \\
\bottomrule
\end{tabular}
\end{table}

La vérification est $\approx 1{,}6$ fois plus rapide que la recherche
originale car elle exécute exactement quatre exponentiations modulaires sur
le candidat connu, sans crible, sans filtre du Théorème~2 ni
pré-sélection Miller–Rabin.

%------------------------------------------------------------------------------
\subsection{Rang dans les records}
%------------------------------------------------------------------------------

Le Tableau~\ref{tab:tousresultats} (Section~\ref{sec:resultats}) liste
l'ensemble des quatre certificats \textsc{PrimeQuest}. L'instance de
$29\,998$ chiffres est actuellement le plus grand certificat
Pocklington–Lehmer inconditionnel dans la famille
$p = 3m(m+1)+1$, $m = 2^a\cdot 3^b - 1$,
obtenu sur un ordinateur portable grand public sans calcul distribué ni
matériel spécialisé.

%==============================================================================
% PRÉAMBULE MINIMAL AUTONOME (décommenter pour compilation stand-alone)
%==============================================================================
% \iffalse   % <-- supprimer cette ligne pour activer le mode autonome
% \documentclass[11pt,a4paper]{article}
% \usepackage[utf8]{inputenc}
% \usepackage[T1]{fontenc}
% \usepackage[french]{babel}
% \usepackage{amsmath,amssymb,amsthm,mathtools,booktabs}
% \usepackage{geometry}\geometry{margin=2.5cm}
% \usepackage[dvipsnames]{xcolor}
% \usepackage{tcolorbox}
% \tcbuselibrary{skins,breakable}
% \newtcolorbox{rigorousbox}[1][]{colback=Green!8,colframe=Green!60!black,
%   arc=1mm,title=\textbf{Rigoureux},fonttitle=\bfseries,breakable,#1}
% \newtcolorbox{resultbox}[1][]{colback=RoyalBlue!5,colframe=RoyalBlue!70!black,
%   arc=1mm,title=\textbf{Résultat calculatoire},fonttitle=\bfseries,breakable,#1}
% \newtheorem{theorem}{Théorème}[section]
% \newtheorem{lemma}[theorem]{Lemme}
% \theoremstyle{definition}
% \newtheorem{definition}[theorem]{Définition}
% \begin{document}
% %==============================================================================
% Certificat de primalité inconditionnelle — premier de 29 998 chiffres
% Famille paramétrique  p = 3m(m+1)+1,  m = 2^a·3^b − 1
% Vérification numérique indépendante — mai 2026
%==============================================================================
%
% Utilisation dans le document principal :
%   \input{certificat_29998_fr}
%
% Compilation autonome :
%   pdflatex certificat_29998_fr.tex
%
%==============================================================================

%------------------------------------------------------------------------------
\section{Certificat inconditionnel pour le premier de 29\,998 chiffres}
\label{sec:cert29998}
%------------------------------------------------------------------------------

Cette section présente le certificat de primalité Pocklington–Lehmer
complet et vérifié indépendamment pour le plus grand premier
actuellement produit par \textsc{PrimeQuest}.

%------------------------------------------------------------------------------
\subsection{Paramètres et structure}
%------------------------------------------------------------------------------

\begin{definition}[Instance de 29\,998 chiffres]\label{def:instance29998}
On pose
\[
  a = 19\,435,\qquad b = 19\,173,
\]
et on définit
\[
  m \;=\; 2^{19435}\cdot 3^{19173} - 1
    \qquad\bigl(14\,999 \text{ chiffres décimaux}\bigr),
\]
\[
  p \;=\; 3m(m+1)+1
    \qquad\bigl(29\,998 \text{ chiffres décimaux}\bigr).
\]
\end{definition}

Le rapport $b/a = 19\,173/19\,435 \approx 0{,}9865$ est voisin de l'optimum
théorique $\rho^* = \log_{10}2/\log_{10}3 \approx 0{,}6309$ minimisant le
nombre de chiffres à $a$ fixé ; la valeur retenue ici cible exactement
$29\,998$ chiffres et a été localisée par l'algorithme \textsc{PrimeQuest v4}
(voir Section~\ref{sec:algo}).

\medskip

\noindent\textbf{Identification.}\quad
Les premiers et derniers chiffres décimaux de $p$ sont :
\begin{equation}\label{eq:p29998chiffres}
  p \;=\; 1602087359509060348500037851549319056715\cdots
         16930852697403293697.
\end{equation}

%------------------------------------------------------------------------------
\subsection{Factorisation de $p-1$}
%------------------------------------------------------------------------------

\begin{rigorousbox}
\begin{lemma}[Factorisation de $p-1$]\label{lem:factorisation29998}
Avec $a, b, m, p$ comme dans la Définition~\ref{def:instance29998},
\[
  p - 1 \;=\; \underbrace{2^{19435}\cdot 3^{19174}}_{F}
              \;\cdot\; m,
\]
où $F := 2^{19435}\cdot 3^{19174}$ possède $15\,000$ chiffres décimaux et
vérifie $F > \sqrt{p}$.
\end{lemma}
\end{rigorousbox}

\begin{proof}
Puisque $m + 1 = 2^{19435}\cdot 3^{19173}$,
\[
  p - 1 \;=\; 3m(m+1) \;=\; 3m\cdot 2^{19435}\cdot 3^{19173}
           \;=\; 2^{19435}\cdot 3^{19174}\cdot m.
\]
Pour la borne $F > \sqrt{p}$, on note que $p < 3(m+1)^2$, donc
$\sqrt{p} < \sqrt{3}\,(m+1) < 3(m+1) = 2^{19435}\cdot 3^{19174} = F$.

\noindent Numériquement, en comparant les logarithmes en base $2$ :
\[
  2\log_2 F
  \;=\; 2\bigl(19435 + 19174\log_2 3\bigr)
  \;\approx\; 99\,650{,}14,
\]
\[
  \log_2 p \;\approx\; 99\,648{,}56,
\]
ce qui confirme $F^2 > p$ avec une marge de $\approx 1{,}58$ bit.
\end{proof}

%------------------------------------------------------------------------------
\subsection{Le certificat Pocklington–Lehmer}
%------------------------------------------------------------------------------

Le Théorème~\ref{thm:Pocklington} (rappelé Section~\ref{sec:family})
requiert, pour chaque premier $q \mid F$ — ici $q \in \{2, 3\}$ —, un entier
$w_q$ vérifiant
\[
  w_q^{p-1} \equiv 1 \pmod{p}
  \qquad\text{et}\qquad
  \gcd\!\bigl(w_q^{(p-1)/q} - 1,\; p\bigr) = 1.
\]

\begin{resultbox}[title={\textbf{Certificat} — premier de 29\,998 chiffres}]
\[
  p \;=\; 3m(m+1)+1,\qquad m = 2^{19435}\cdot 3^{19173}-1.
\]

\medskip
\noindent\textbf{Partie connue de $p-1$ :}
$\;F = 2^{19435}\cdot 3^{19174},\quad F > \sqrt{p}.\quad$\checkmark

\medskip
\renewcommand{\arraystretch}{1.4}
\begin{tabular}{@{}c l l@{}}
\toprule
Premier $q$ & Témoin $w_q$ & Conditions vérifiées \\
\midrule
$2$ & $w_2 = 5$ &
  $5^{p-1}\equiv 1\pmod{p}$\quad et\quad
  $\gcd(5^{(p-1)/2}-1,\,p)=1$ \\
$3$ & $w_3 = 7$ &
  $7^{p-1}\equiv 1\pmod{p}$\quad et\quad
  $\gcd(7^{(p-1)/3}-1,\,p)=1$ \\
\bottomrule
\end{tabular}

\medskip
\noindent\textbf{Conclusion :} $p$ est premier (preuve déterministe). \checkmark
\end{resultbox}

%------------------------------------------------------------------------------
\subsection{Vérification numérique indépendante}
%------------------------------------------------------------------------------

Les quatre conditions du certificat ont été vérifiées de façon indépendante
à l'aide de \texttt{Python 3.12} et de son arithmétique entière en précision
arbitraire (backend GMP). Le script ne comporte aucune étape probabiliste ;
toutes les opérations sont des arithmétiques modulaires exactes.

\begin{rigorousbox}[title={\textbf{Attestation de vérification} — calcul indépendant}]
\textbf{Environnement :} Python~3.12,
\texttt{sys.set\_int\_max\_str\_digits(40000)} ;
exponentiations modulaires via \texttt{pow(base, exp, mod)}.\\[4pt]

\textbf{Étape 1 — Reconstruction de $p$.}\quad
Calcul de $m = 2^{19435}\cdot 3^{19173}-1$ (14\,999 chiffres) et de
$p = 3m(m+1)+1$ (29\,998 chiffres) à partir des paramètres bruts.
Premiers 40 chiffres de $p$ :
\texttt{1602087359509060348500037851549319056715} ;
20 derniers : \texttt{16930852697403293697}.
Correspond exactement au certificat stocké.\\[4pt]

\textbf{Étape 2 — Borne Pocklington.}
\[
  2\log_2 F = 99\,650{,}1420 \;>\; \log_2 p = 99\,648{,}5570.\quad\checkmark
\]

\textbf{Étape 3 — Témoin $w_2 = 5$ pour $q = 2$.}
\begin{align*}
  5^{p-1} &\equiv 1 \pmod{p},\quad\checkmark \\
  \gcd(5^{(p-1)/2}-1,\,p) &= 1.\quad\checkmark
\end{align*}

\textbf{Étape 4 — Témoin $w_3 = 7$ pour $q = 3$.}
\begin{align*}
  7^{p-1} &\equiv 1 \pmod{p},\quad\checkmark \\
  \gcd(7^{(p-1)/3}-1,\,p) &= 1.\quad\checkmark
\end{align*}

\textbf{Durée totale :} $6\,945$\,s $\approx 1$\,h\,$56$\,min
(cœur unique, Python GMP, sans parallélisme).\\[4pt]

\textbf{Conclusion :} les quatre conditions Pocklington–Lehmer sont satisfaites.
Le nombre $p$ est \emph{premier de façon inconditionnelle}.
\end{rigorousbox}

%------------------------------------------------------------------------------
\subsection{Comparaison avec le calcul original}
%------------------------------------------------------------------------------

\begin{table}[h]
\centering
\caption{Recherche originale vs.\ vérification indépendante du certificat
         pour le premier de 29\,998 chiffres.}
\label{tab:comparaison29998}
\renewcommand{\arraystretch}{1.3}
\begin{tabular}{@{}lll@{}}
\toprule
 & \textbf{PrimeQuest v4 (recherche)} & \textbf{Vérification indépendante} \\
\midrule
Objectif       & Trouver un candidat premier         & Vérifier le certificat \\
Machine        & ARM64, 9 cœurs (Snapdragon X)       & Cœur unique (Python GMP) \\
Durée          & $11\,146$\,s ($3$\,h\,$06$\,min)   & $6\,945$\,s ($1$\,h\,$56$\,min) \\
Tests MR       & $286$ (pré-filtre probabiliste)     & — (aucune étape probabiliste) \\
Résultat       & Certificat produit                  & Certificat confirmé \\
\bottomrule
\end{tabular}
\end{table}

La vérification est $\approx 1{,}6$ fois plus rapide que la recherche
originale car elle exécute exactement quatre exponentiations modulaires sur
le candidat connu, sans crible, sans filtre du Théorème~2 ni
pré-sélection Miller–Rabin.

%------------------------------------------------------------------------------
\subsection{Rang dans les records}
%------------------------------------------------------------------------------

Le Tableau~\ref{tab:tousresultats} (Section~\ref{sec:resultats}) liste
l'ensemble des quatre certificats \textsc{PrimeQuest}. L'instance de
$29\,998$ chiffres est actuellement le plus grand certificat
Pocklington–Lehmer inconditionnel dans la famille
$p = 3m(m+1)+1$, $m = 2^a\cdot 3^b - 1$,
obtenu sur un ordinateur portable grand public sans calcul distribué ni
matériel spécialisé.

%==============================================================================
% PRÉAMBULE MINIMAL AUTONOME (décommenter pour compilation stand-alone)
%==============================================================================
% \iffalse   % <-- supprimer cette ligne pour activer le mode autonome
% \documentclass[11pt,a4paper]{article}
% \usepackage[utf8]{inputenc}
% \usepackage[T1]{fontenc}
% \usepackage[french]{babel}
% \usepackage{amsmath,amssymb,amsthm,mathtools,booktabs}
% \usepackage{geometry}\geometry{margin=2.5cm}
% \usepackage[dvipsnames]{xcolor}
% \usepackage{tcolorbox}
% \tcbuselibrary{skins,breakable}
% \newtcolorbox{rigorousbox}[1][]{colback=Green!8,colframe=Green!60!black,
%   arc=1mm,title=\textbf{Rigoureux},fonttitle=\bfseries,breakable,#1}
% \newtcolorbox{resultbox}[1][]{colback=RoyalBlue!5,colframe=RoyalBlue!70!black,
%   arc=1mm,title=\textbf{Résultat calculatoire},fonttitle=\bfseries,breakable,#1}
% \newtheorem{theorem}{Théorème}[section]
% \newtheorem{lemma}[theorem]{Lemme}
% \theoremstyle{definition}
% \newtheorem{definition}[theorem]{Définition}
% \begin{document}
% \input{certificat_29998_fr}
% \end{document}
% \fi

% \end{document}
% \fi

%
% Compilation autonome :
%   pdflatex certificat_29998_fr.tex
%
%==============================================================================

%------------------------------------------------------------------------------
\section{Certificat inconditionnel pour le premier de 29\,998 chiffres}
\label{sec:cert29998}
%------------------------------------------------------------------------------

Cette section présente le certificat de primalité Pocklington–Lehmer
complet et vérifié indépendamment pour le plus grand premier
actuellement produit par \textsc{PrimeQuest}.

%------------------------------------------------------------------------------
\subsection{Paramètres et structure}
%------------------------------------------------------------------------------

\begin{definition}[Instance de 29\,998 chiffres]\label{def:instance29998}
On pose
\[
  a = 19\,435,\qquad b = 19\,173,
\]
et on définit
\[
  m \;=\; 2^{19435}\cdot 3^{19173} - 1
    \qquad\bigl(14\,999 \text{ chiffres décimaux}\bigr),
\]
\[
  p \;=\; 3m(m+1)+1
    \qquad\bigl(29\,998 \text{ chiffres décimaux}\bigr).
\]
\end{definition}

Le rapport $b/a = 19\,173/19\,435 \approx 0{,}9865$ est voisin de l'optimum
théorique $\rho^* = \log_{10}2/\log_{10}3 \approx 0{,}6309$ minimisant le
nombre de chiffres à $a$ fixé ; la valeur retenue ici cible exactement
$29\,998$ chiffres et a été localisée par l'algorithme \textsc{PrimeQuest v4}
(voir Section~\ref{sec:algo}).

\medskip

\noindent\textbf{Identification.}\quad
Les premiers et derniers chiffres décimaux de $p$ sont :
\begin{equation}\label{eq:p29998chiffres}
  p \;=\; 1602087359509060348500037851549319056715\cdots
         16930852697403293697.
\end{equation}

%------------------------------------------------------------------------------
\subsection{Factorisation de $p-1$}
%------------------------------------------------------------------------------

\begin{rigorousbox}
\begin{lemma}[Factorisation de $p-1$]\label{lem:factorisation29998}
Avec $a, b, m, p$ comme dans la Définition~\ref{def:instance29998},
\[
  p - 1 \;=\; \underbrace{2^{19435}\cdot 3^{19174}}_{F}
              \;\cdot\; m,
\]
où $F := 2^{19435}\cdot 3^{19174}$ possède $15\,000$ chiffres décimaux et
vérifie $F > \sqrt{p}$.
\end{lemma}
\end{rigorousbox}

\begin{proof}
Puisque $m + 1 = 2^{19435}\cdot 3^{19173}$,
\[
  p - 1 \;=\; 3m(m+1) \;=\; 3m\cdot 2^{19435}\cdot 3^{19173}
           \;=\; 2^{19435}\cdot 3^{19174}\cdot m.
\]
Pour la borne $F > \sqrt{p}$, on note que $p < 3(m+1)^2$, donc
$\sqrt{p} < \sqrt{3}\,(m+1) < 3(m+1) = 2^{19435}\cdot 3^{19174} = F$.

\noindent Numériquement, en comparant les logarithmes en base $2$ :
\[
  2\log_2 F
  \;=\; 2\bigl(19435 + 19174\log_2 3\bigr)
  \;\approx\; 99\,650{,}14,
\]
\[
  \log_2 p \;\approx\; 99\,648{,}56,
\]
ce qui confirme $F^2 > p$ avec une marge de $\approx 1{,}58$ bit.
\end{proof}

%------------------------------------------------------------------------------
\subsection{Le certificat Pocklington–Lehmer}
%------------------------------------------------------------------------------

Le Théorème~\ref{thm:Pocklington} (rappelé Section~\ref{sec:family})
requiert, pour chaque premier $q \mid F$ — ici $q \in \{2, 3\}$ —, un entier
$w_q$ vérifiant
\[
  w_q^{p-1} \equiv 1 \pmod{p}
  \qquad\text{et}\qquad
  \gcd\!\bigl(w_q^{(p-1)/q} - 1,\; p\bigr) = 1.
\]

\begin{resultbox}[title={\textbf{Certificat} — premier de 29\,998 chiffres}]
\[
  p \;=\; 3m(m+1)+1,\qquad m = 2^{19435}\cdot 3^{19173}-1.
\]

\medskip
\noindent\textbf{Partie connue de $p-1$ :}
$\;F = 2^{19435}\cdot 3^{19174},\quad F > \sqrt{p}.\quad$\checkmark

\medskip
\renewcommand{\arraystretch}{1.4}
\begin{tabular}{@{}c l l@{}}
\toprule
Premier $q$ & Témoin $w_q$ & Conditions vérifiées \\
\midrule
$2$ & $w_2 = 5$ &
  $5^{p-1}\equiv 1\pmod{p}$\quad et\quad
  $\gcd(5^{(p-1)/2}-1,\,p)=1$ \\
$3$ & $w_3 = 7$ &
  $7^{p-1}\equiv 1\pmod{p}$\quad et\quad
  $\gcd(7^{(p-1)/3}-1,\,p)=1$ \\
\bottomrule
\end{tabular}

\medskip
\noindent\textbf{Conclusion :} $p$ est premier (preuve déterministe). \checkmark
\end{resultbox}

%------------------------------------------------------------------------------
\subsection{Vérification numérique indépendante}
%------------------------------------------------------------------------------

Les quatre conditions du certificat ont été vérifiées de façon indépendante
à l'aide de \texttt{Python 3.12} et de son arithmétique entière en précision
arbitraire (backend GMP). Le script ne comporte aucune étape probabiliste ;
toutes les opérations sont des arithmétiques modulaires exactes.

\begin{rigorousbox}[title={\textbf{Attestation de vérification} — calcul indépendant}]
\textbf{Environnement :} Python~3.12,
\texttt{sys.set\_int\_max\_str\_digits(40000)} ;
exponentiations modulaires via \texttt{pow(base, exp, mod)}.\\[4pt]

\textbf{Étape 1 — Reconstruction de $p$.}\quad
Calcul de $m = 2^{19435}\cdot 3^{19173}-1$ (14\,999 chiffres) et de
$p = 3m(m+1)+1$ (29\,998 chiffres) à partir des paramètres bruts.
Premiers 40 chiffres de $p$ :
\texttt{1602087359509060348500037851549319056715} ;
20 derniers : \texttt{16930852697403293697}.
Correspond exactement au certificat stocké.\\[4pt]

\textbf{Étape 2 — Borne Pocklington.}
\[
  2\log_2 F = 99\,650{,}1420 \;>\; \log_2 p = 99\,648{,}5570.\quad\checkmark
\]

\textbf{Étape 3 — Témoin $w_2 = 5$ pour $q = 2$.}
\begin{align*}
  5^{p-1} &\equiv 1 \pmod{p},\quad\checkmark \\
  \gcd(5^{(p-1)/2}-1,\,p) &= 1.\quad\checkmark
\end{align*}

\textbf{Étape 4 — Témoin $w_3 = 7$ pour $q = 3$.}
\begin{align*}
  7^{p-1} &\equiv 1 \pmod{p},\quad\checkmark \\
  \gcd(7^{(p-1)/3}-1,\,p) &= 1.\quad\checkmark
\end{align*}

\textbf{Durée totale :} $6\,945$\,s $\approx 1$\,h\,$56$\,min
(cœur unique, Python GMP, sans parallélisme).\\[4pt]

\textbf{Conclusion :} les quatre conditions Pocklington–Lehmer sont satisfaites.
Le nombre $p$ est \emph{premier de façon inconditionnelle}.
\end{rigorousbox}

%------------------------------------------------------------------------------
\subsection{Comparaison avec le calcul original}
%------------------------------------------------------------------------------

\begin{table}[h]
\centering
\caption{Recherche originale vs.\ vérification indépendante du certificat
         pour le premier de 29\,998 chiffres.}
\label{tab:comparaison29998}
\renewcommand{\arraystretch}{1.3}
\begin{tabular}{@{}lll@{}}
\toprule
 & \textbf{PrimeQuest v4 (recherche)} & \textbf{Vérification indépendante} \\
\midrule
Objectif       & Trouver un candidat premier         & Vérifier le certificat \\
Machine        & ARM64, 9 cœurs (Snapdragon X)       & Cœur unique (Python GMP) \\
Durée          & $11\,146$\,s ($3$\,h\,$06$\,min)   & $6\,945$\,s ($1$\,h\,$56$\,min) \\
Tests MR       & $286$ (pré-filtre probabiliste)     & — (aucune étape probabiliste) \\
Résultat       & Certificat produit                  & Certificat confirmé \\
\bottomrule
\end{tabular}
\end{table}

La vérification est $\approx 1{,}6$ fois plus rapide que la recherche
originale car elle exécute exactement quatre exponentiations modulaires sur
le candidat connu, sans crible, sans filtre du Théorème~2 ni
pré-sélection Miller–Rabin.

%------------------------------------------------------------------------------
\subsection{Rang dans les records}
%------------------------------------------------------------------------------

Le Tableau~\ref{tab:tousresultats} (Section~\ref{sec:resultats}) liste
l'ensemble des quatre certificats \textsc{PrimeQuest}. L'instance de
$29\,998$ chiffres est actuellement le plus grand certificat
Pocklington–Lehmer inconditionnel dans la famille
$p = 3m(m+1)+1$, $m = 2^a\cdot 3^b - 1$,
obtenu sur un ordinateur portable grand public sans calcul distribué ni
matériel spécialisé.

%==============================================================================
% PRÉAMBULE MINIMAL AUTONOME (décommenter pour compilation stand-alone)
%==============================================================================
% \iffalse   % <-- supprimer cette ligne pour activer le mode autonome
% \documentclass[11pt,a4paper]{article}
% \usepackage[utf8]{inputenc}
% \usepackage[T1]{fontenc}
% \usepackage[french]{babel}
% \usepackage{amsmath,amssymb,amsthm,mathtools,booktabs}
% \usepackage{geometry}\geometry{margin=2.5cm}
% \usepackage[dvipsnames]{xcolor}
% \usepackage{tcolorbox}
% \tcbuselibrary{skins,breakable}
% \newtcolorbox{rigorousbox}[1][]{colback=Green!8,colframe=Green!60!black,
%   arc=1mm,title=\textbf{Rigoureux},fonttitle=\bfseries,breakable,#1}
% \newtcolorbox{resultbox}[1][]{colback=RoyalBlue!5,colframe=RoyalBlue!70!black,
%   arc=1mm,title=\textbf{Résultat calculatoire},fonttitle=\bfseries,breakable,#1}
% \newtheorem{theorem}{Théorème}[section]
% \newtheorem{lemma}[theorem]{Lemme}
% \theoremstyle{definition}
% \newtheorem{definition}[theorem]{Définition}
% \begin{document}
% %==============================================================================
% Certificat de primalité inconditionnelle — premier de 29 998 chiffres
% Famille paramétrique  p = 3m(m+1)+1,  m = 2^a·3^b − 1
% Vérification numérique indépendante — mai 2026
%==============================================================================
%
% Utilisation dans le document principal :
%   %==============================================================================
% Certificat de primalité inconditionnelle — premier de 29 998 chiffres
% Famille paramétrique  p = 3m(m+1)+1,  m = 2^a·3^b − 1
% Vérification numérique indépendante — mai 2026
%==============================================================================
%
% Utilisation dans le document principal :
%   \input{certificat_29998_fr}
%
% Compilation autonome :
%   pdflatex certificat_29998_fr.tex
%
%==============================================================================

%------------------------------------------------------------------------------
\section{Certificat inconditionnel pour le premier de 29\,998 chiffres}
\label{sec:cert29998}
%------------------------------------------------------------------------------

Cette section présente le certificat de primalité Pocklington–Lehmer
complet et vérifié indépendamment pour le plus grand premier
actuellement produit par \textsc{PrimeQuest}.

%------------------------------------------------------------------------------
\subsection{Paramètres et structure}
%------------------------------------------------------------------------------

\begin{definition}[Instance de 29\,998 chiffres]\label{def:instance29998}
On pose
\[
  a = 19\,435,\qquad b = 19\,173,
\]
et on définit
\[
  m \;=\; 2^{19435}\cdot 3^{19173} - 1
    \qquad\bigl(14\,999 \text{ chiffres décimaux}\bigr),
\]
\[
  p \;=\; 3m(m+1)+1
    \qquad\bigl(29\,998 \text{ chiffres décimaux}\bigr).
\]
\end{definition}

Le rapport $b/a = 19\,173/19\,435 \approx 0{,}9865$ est voisin de l'optimum
théorique $\rho^* = \log_{10}2/\log_{10}3 \approx 0{,}6309$ minimisant le
nombre de chiffres à $a$ fixé ; la valeur retenue ici cible exactement
$29\,998$ chiffres et a été localisée par l'algorithme \textsc{PrimeQuest v4}
(voir Section~\ref{sec:algo}).

\medskip

\noindent\textbf{Identification.}\quad
Les premiers et derniers chiffres décimaux de $p$ sont :
\begin{equation}\label{eq:p29998chiffres}
  p \;=\; 1602087359509060348500037851549319056715\cdots
         16930852697403293697.
\end{equation}

%------------------------------------------------------------------------------
\subsection{Factorisation de $p-1$}
%------------------------------------------------------------------------------

\begin{rigorousbox}
\begin{lemma}[Factorisation de $p-1$]\label{lem:factorisation29998}
Avec $a, b, m, p$ comme dans la Définition~\ref{def:instance29998},
\[
  p - 1 \;=\; \underbrace{2^{19435}\cdot 3^{19174}}_{F}
              \;\cdot\; m,
\]
où $F := 2^{19435}\cdot 3^{19174}$ possède $15\,000$ chiffres décimaux et
vérifie $F > \sqrt{p}$.
\end{lemma}
\end{rigorousbox}

\begin{proof}
Puisque $m + 1 = 2^{19435}\cdot 3^{19173}$,
\[
  p - 1 \;=\; 3m(m+1) \;=\; 3m\cdot 2^{19435}\cdot 3^{19173}
           \;=\; 2^{19435}\cdot 3^{19174}\cdot m.
\]
Pour la borne $F > \sqrt{p}$, on note que $p < 3(m+1)^2$, donc
$\sqrt{p} < \sqrt{3}\,(m+1) < 3(m+1) = 2^{19435}\cdot 3^{19174} = F$.

\noindent Numériquement, en comparant les logarithmes en base $2$ :
\[
  2\log_2 F
  \;=\; 2\bigl(19435 + 19174\log_2 3\bigr)
  \;\approx\; 99\,650{,}14,
\]
\[
  \log_2 p \;\approx\; 99\,648{,}56,
\]
ce qui confirme $F^2 > p$ avec une marge de $\approx 1{,}58$ bit.
\end{proof}

%------------------------------------------------------------------------------
\subsection{Le certificat Pocklington–Lehmer}
%------------------------------------------------------------------------------

Le Théorème~\ref{thm:Pocklington} (rappelé Section~\ref{sec:family})
requiert, pour chaque premier $q \mid F$ — ici $q \in \{2, 3\}$ —, un entier
$w_q$ vérifiant
\[
  w_q^{p-1} \equiv 1 \pmod{p}
  \qquad\text{et}\qquad
  \gcd\!\bigl(w_q^{(p-1)/q} - 1,\; p\bigr) = 1.
\]

\begin{resultbox}[title={\textbf{Certificat} — premier de 29\,998 chiffres}]
\[
  p \;=\; 3m(m+1)+1,\qquad m = 2^{19435}\cdot 3^{19173}-1.
\]

\medskip
\noindent\textbf{Partie connue de $p-1$ :}
$\;F = 2^{19435}\cdot 3^{19174},\quad F > \sqrt{p}.\quad$\checkmark

\medskip
\renewcommand{\arraystretch}{1.4}
\begin{tabular}{@{}c l l@{}}
\toprule
Premier $q$ & Témoin $w_q$ & Conditions vérifiées \\
\midrule
$2$ & $w_2 = 5$ &
  $5^{p-1}\equiv 1\pmod{p}$\quad et\quad
  $\gcd(5^{(p-1)/2}-1,\,p)=1$ \\
$3$ & $w_3 = 7$ &
  $7^{p-1}\equiv 1\pmod{p}$\quad et\quad
  $\gcd(7^{(p-1)/3}-1,\,p)=1$ \\
\bottomrule
\end{tabular}

\medskip
\noindent\textbf{Conclusion :} $p$ est premier (preuve déterministe). \checkmark
\end{resultbox}

%------------------------------------------------------------------------------
\subsection{Vérification numérique indépendante}
%------------------------------------------------------------------------------

Les quatre conditions du certificat ont été vérifiées de façon indépendante
à l'aide de \texttt{Python 3.12} et de son arithmétique entière en précision
arbitraire (backend GMP). Le script ne comporte aucune étape probabiliste ;
toutes les opérations sont des arithmétiques modulaires exactes.

\begin{rigorousbox}[title={\textbf{Attestation de vérification} — calcul indépendant}]
\textbf{Environnement :} Python~3.12,
\texttt{sys.set\_int\_max\_str\_digits(40000)} ;
exponentiations modulaires via \texttt{pow(base, exp, mod)}.\\[4pt]

\textbf{Étape 1 — Reconstruction de $p$.}\quad
Calcul de $m = 2^{19435}\cdot 3^{19173}-1$ (14\,999 chiffres) et de
$p = 3m(m+1)+1$ (29\,998 chiffres) à partir des paramètres bruts.
Premiers 40 chiffres de $p$ :
\texttt{1602087359509060348500037851549319056715} ;
20 derniers : \texttt{16930852697403293697}.
Correspond exactement au certificat stocké.\\[4pt]

\textbf{Étape 2 — Borne Pocklington.}
\[
  2\log_2 F = 99\,650{,}1420 \;>\; \log_2 p = 99\,648{,}5570.\quad\checkmark
\]

\textbf{Étape 3 — Témoin $w_2 = 5$ pour $q = 2$.}
\begin{align*}
  5^{p-1} &\equiv 1 \pmod{p},\quad\checkmark \\
  \gcd(5^{(p-1)/2}-1,\,p) &= 1.\quad\checkmark
\end{align*}

\textbf{Étape 4 — Témoin $w_3 = 7$ pour $q = 3$.}
\begin{align*}
  7^{p-1} &\equiv 1 \pmod{p},\quad\checkmark \\
  \gcd(7^{(p-1)/3}-1,\,p) &= 1.\quad\checkmark
\end{align*}

\textbf{Durée totale :} $6\,945$\,s $\approx 1$\,h\,$56$\,min
(cœur unique, Python GMP, sans parallélisme).\\[4pt]

\textbf{Conclusion :} les quatre conditions Pocklington–Lehmer sont satisfaites.
Le nombre $p$ est \emph{premier de façon inconditionnelle}.
\end{rigorousbox}

%------------------------------------------------------------------------------
\subsection{Comparaison avec le calcul original}
%------------------------------------------------------------------------------

\begin{table}[h]
\centering
\caption{Recherche originale vs.\ vérification indépendante du certificat
         pour le premier de 29\,998 chiffres.}
\label{tab:comparaison29998}
\renewcommand{\arraystretch}{1.3}
\begin{tabular}{@{}lll@{}}
\toprule
 & \textbf{PrimeQuest v4 (recherche)} & \textbf{Vérification indépendante} \\
\midrule
Objectif       & Trouver un candidat premier         & Vérifier le certificat \\
Machine        & ARM64, 9 cœurs (Snapdragon X)       & Cœur unique (Python GMP) \\
Durée          & $11\,146$\,s ($3$\,h\,$06$\,min)   & $6\,945$\,s ($1$\,h\,$56$\,min) \\
Tests MR       & $286$ (pré-filtre probabiliste)     & — (aucune étape probabiliste) \\
Résultat       & Certificat produit                  & Certificat confirmé \\
\bottomrule
\end{tabular}
\end{table}

La vérification est $\approx 1{,}6$ fois plus rapide que la recherche
originale car elle exécute exactement quatre exponentiations modulaires sur
le candidat connu, sans crible, sans filtre du Théorème~2 ni
pré-sélection Miller–Rabin.

%------------------------------------------------------------------------------
\subsection{Rang dans les records}
%------------------------------------------------------------------------------

Le Tableau~\ref{tab:tousresultats} (Section~\ref{sec:resultats}) liste
l'ensemble des quatre certificats \textsc{PrimeQuest}. L'instance de
$29\,998$ chiffres est actuellement le plus grand certificat
Pocklington–Lehmer inconditionnel dans la famille
$p = 3m(m+1)+1$, $m = 2^a\cdot 3^b - 1$,
obtenu sur un ordinateur portable grand public sans calcul distribué ni
matériel spécialisé.

%==============================================================================
% PRÉAMBULE MINIMAL AUTONOME (décommenter pour compilation stand-alone)
%==============================================================================
% \iffalse   % <-- supprimer cette ligne pour activer le mode autonome
% \documentclass[11pt,a4paper]{article}
% \usepackage[utf8]{inputenc}
% \usepackage[T1]{fontenc}
% \usepackage[french]{babel}
% \usepackage{amsmath,amssymb,amsthm,mathtools,booktabs}
% \usepackage{geometry}\geometry{margin=2.5cm}
% \usepackage[dvipsnames]{xcolor}
% \usepackage{tcolorbox}
% \tcbuselibrary{skins,breakable}
% \newtcolorbox{rigorousbox}[1][]{colback=Green!8,colframe=Green!60!black,
%   arc=1mm,title=\textbf{Rigoureux},fonttitle=\bfseries,breakable,#1}
% \newtcolorbox{resultbox}[1][]{colback=RoyalBlue!5,colframe=RoyalBlue!70!black,
%   arc=1mm,title=\textbf{Résultat calculatoire},fonttitle=\bfseries,breakable,#1}
% \newtheorem{theorem}{Théorème}[section]
% \newtheorem{lemma}[theorem]{Lemme}
% \theoremstyle{definition}
% \newtheorem{definition}[theorem]{Définition}
% \begin{document}
% \input{certificat_29998_fr}
% \end{document}
% \fi

%
% Compilation autonome :
%   pdflatex certificat_29998_fr.tex
%
%==============================================================================

%------------------------------------------------------------------------------
\section{Certificat inconditionnel pour le premier de 29\,998 chiffres}
\label{sec:cert29998}
%------------------------------------------------------------------------------

Cette section présente le certificat de primalité Pocklington–Lehmer
complet et vérifié indépendamment pour le plus grand premier
actuellement produit par \textsc{PrimeQuest}.

%------------------------------------------------------------------------------
\subsection{Paramètres et structure}
%------------------------------------------------------------------------------

\begin{definition}[Instance de 29\,998 chiffres]\label{def:instance29998}
On pose
\[
  a = 19\,435,\qquad b = 19\,173,
\]
et on définit
\[
  m \;=\; 2^{19435}\cdot 3^{19173} - 1
    \qquad\bigl(14\,999 \text{ chiffres décimaux}\bigr),
\]
\[
  p \;=\; 3m(m+1)+1
    \qquad\bigl(29\,998 \text{ chiffres décimaux}\bigr).
\]
\end{definition}

Le rapport $b/a = 19\,173/19\,435 \approx 0{,}9865$ est voisin de l'optimum
théorique $\rho^* = \log_{10}2/\log_{10}3 \approx 0{,}6309$ minimisant le
nombre de chiffres à $a$ fixé ; la valeur retenue ici cible exactement
$29\,998$ chiffres et a été localisée par l'algorithme \textsc{PrimeQuest v4}
(voir Section~\ref{sec:algo}).

\medskip

\noindent\textbf{Identification.}\quad
Les premiers et derniers chiffres décimaux de $p$ sont :
\begin{equation}\label{eq:p29998chiffres}
  p \;=\; 1602087359509060348500037851549319056715\cdots
         16930852697403293697.
\end{equation}

%------------------------------------------------------------------------------
\subsection{Factorisation de $p-1$}
%------------------------------------------------------------------------------

\begin{rigorousbox}
\begin{lemma}[Factorisation de $p-1$]\label{lem:factorisation29998}
Avec $a, b, m, p$ comme dans la Définition~\ref{def:instance29998},
\[
  p - 1 \;=\; \underbrace{2^{19435}\cdot 3^{19174}}_{F}
              \;\cdot\; m,
\]
où $F := 2^{19435}\cdot 3^{19174}$ possède $15\,000$ chiffres décimaux et
vérifie $F > \sqrt{p}$.
\end{lemma}
\end{rigorousbox}

\begin{proof}
Puisque $m + 1 = 2^{19435}\cdot 3^{19173}$,
\[
  p - 1 \;=\; 3m(m+1) \;=\; 3m\cdot 2^{19435}\cdot 3^{19173}
           \;=\; 2^{19435}\cdot 3^{19174}\cdot m.
\]
Pour la borne $F > \sqrt{p}$, on note que $p < 3(m+1)^2$, donc
$\sqrt{p} < \sqrt{3}\,(m+1) < 3(m+1) = 2^{19435}\cdot 3^{19174} = F$.

\noindent Numériquement, en comparant les logarithmes en base $2$ :
\[
  2\log_2 F
  \;=\; 2\bigl(19435 + 19174\log_2 3\bigr)
  \;\approx\; 99\,650{,}14,
\]
\[
  \log_2 p \;\approx\; 99\,648{,}56,
\]
ce qui confirme $F^2 > p$ avec une marge de $\approx 1{,}58$ bit.
\end{proof}

%------------------------------------------------------------------------------
\subsection{Le certificat Pocklington–Lehmer}
%------------------------------------------------------------------------------

Le Théorème~\ref{thm:Pocklington} (rappelé Section~\ref{sec:family})
requiert, pour chaque premier $q \mid F$ — ici $q \in \{2, 3\}$ —, un entier
$w_q$ vérifiant
\[
  w_q^{p-1} \equiv 1 \pmod{p}
  \qquad\text{et}\qquad
  \gcd\!\bigl(w_q^{(p-1)/q} - 1,\; p\bigr) = 1.
\]

\begin{resultbox}[title={\textbf{Certificat} — premier de 29\,998 chiffres}]
\[
  p \;=\; 3m(m+1)+1,\qquad m = 2^{19435}\cdot 3^{19173}-1.
\]

\medskip
\noindent\textbf{Partie connue de $p-1$ :}
$\;F = 2^{19435}\cdot 3^{19174},\quad F > \sqrt{p}.\quad$\checkmark

\medskip
\renewcommand{\arraystretch}{1.4}
\begin{tabular}{@{}c l l@{}}
\toprule
Premier $q$ & Témoin $w_q$ & Conditions vérifiées \\
\midrule
$2$ & $w_2 = 5$ &
  $5^{p-1}\equiv 1\pmod{p}$\quad et\quad
  $\gcd(5^{(p-1)/2}-1,\,p)=1$ \\
$3$ & $w_3 = 7$ &
  $7^{p-1}\equiv 1\pmod{p}$\quad et\quad
  $\gcd(7^{(p-1)/3}-1,\,p)=1$ \\
\bottomrule
\end{tabular}

\medskip
\noindent\textbf{Conclusion :} $p$ est premier (preuve déterministe). \checkmark
\end{resultbox}

%------------------------------------------------------------------------------
\subsection{Vérification numérique indépendante}
%------------------------------------------------------------------------------

Les quatre conditions du certificat ont été vérifiées de façon indépendante
à l'aide de \texttt{Python 3.12} et de son arithmétique entière en précision
arbitraire (backend GMP). Le script ne comporte aucune étape probabiliste ;
toutes les opérations sont des arithmétiques modulaires exactes.

\begin{rigorousbox}[title={\textbf{Attestation de vérification} — calcul indépendant}]
\textbf{Environnement :} Python~3.12,
\texttt{sys.set\_int\_max\_str\_digits(40000)} ;
exponentiations modulaires via \texttt{pow(base, exp, mod)}.\\[4pt]

\textbf{Étape 1 — Reconstruction de $p$.}\quad
Calcul de $m = 2^{19435}\cdot 3^{19173}-1$ (14\,999 chiffres) et de
$p = 3m(m+1)+1$ (29\,998 chiffres) à partir des paramètres bruts.
Premiers 40 chiffres de $p$ :
\texttt{1602087359509060348500037851549319056715} ;
20 derniers : \texttt{16930852697403293697}.
Correspond exactement au certificat stocké.\\[4pt]

\textbf{Étape 2 — Borne Pocklington.}
\[
  2\log_2 F = 99\,650{,}1420 \;>\; \log_2 p = 99\,648{,}5570.\quad\checkmark
\]

\textbf{Étape 3 — Témoin $w_2 = 5$ pour $q = 2$.}
\begin{align*}
  5^{p-1} &\equiv 1 \pmod{p},\quad\checkmark \\
  \gcd(5^{(p-1)/2}-1,\,p) &= 1.\quad\checkmark
\end{align*}

\textbf{Étape 4 — Témoin $w_3 = 7$ pour $q = 3$.}
\begin{align*}
  7^{p-1} &\equiv 1 \pmod{p},\quad\checkmark \\
  \gcd(7^{(p-1)/3}-1,\,p) &= 1.\quad\checkmark
\end{align*}

\textbf{Durée totale :} $6\,945$\,s $\approx 1$\,h\,$56$\,min
(cœur unique, Python GMP, sans parallélisme).\\[4pt]

\textbf{Conclusion :} les quatre conditions Pocklington–Lehmer sont satisfaites.
Le nombre $p$ est \emph{premier de façon inconditionnelle}.
\end{rigorousbox}

%------------------------------------------------------------------------------
\subsection{Comparaison avec le calcul original}
%------------------------------------------------------------------------------

\begin{table}[h]
\centering
\caption{Recherche originale vs.\ vérification indépendante du certificat
         pour le premier de 29\,998 chiffres.}
\label{tab:comparaison29998}
\renewcommand{\arraystretch}{1.3}
\begin{tabular}{@{}lll@{}}
\toprule
 & \textbf{PrimeQuest v4 (recherche)} & \textbf{Vérification indépendante} \\
\midrule
Objectif       & Trouver un candidat premier         & Vérifier le certificat \\
Machine        & ARM64, 9 cœurs (Snapdragon X)       & Cœur unique (Python GMP) \\
Durée          & $11\,146$\,s ($3$\,h\,$06$\,min)   & $6\,945$\,s ($1$\,h\,$56$\,min) \\
Tests MR       & $286$ (pré-filtre probabiliste)     & — (aucune étape probabiliste) \\
Résultat       & Certificat produit                  & Certificat confirmé \\
\bottomrule
\end{tabular}
\end{table}

La vérification est $\approx 1{,}6$ fois plus rapide que la recherche
originale car elle exécute exactement quatre exponentiations modulaires sur
le candidat connu, sans crible, sans filtre du Théorème~2 ni
pré-sélection Miller–Rabin.

%------------------------------------------------------------------------------
\subsection{Rang dans les records}
%------------------------------------------------------------------------------

Le Tableau~\ref{tab:tousresultats} (Section~\ref{sec:resultats}) liste
l'ensemble des quatre certificats \textsc{PrimeQuest}. L'instance de
$29\,998$ chiffres est actuellement le plus grand certificat
Pocklington–Lehmer inconditionnel dans la famille
$p = 3m(m+1)+1$, $m = 2^a\cdot 3^b - 1$,
obtenu sur un ordinateur portable grand public sans calcul distribué ni
matériel spécialisé.

%==============================================================================
% PRÉAMBULE MINIMAL AUTONOME (décommenter pour compilation stand-alone)
%==============================================================================
% \iffalse   % <-- supprimer cette ligne pour activer le mode autonome
% \documentclass[11pt,a4paper]{article}
% \usepackage[utf8]{inputenc}
% \usepackage[T1]{fontenc}
% \usepackage[french]{babel}
% \usepackage{amsmath,amssymb,amsthm,mathtools,booktabs}
% \usepackage{geometry}\geometry{margin=2.5cm}
% \usepackage[dvipsnames]{xcolor}
% \usepackage{tcolorbox}
% \tcbuselibrary{skins,breakable}
% \newtcolorbox{rigorousbox}[1][]{colback=Green!8,colframe=Green!60!black,
%   arc=1mm,title=\textbf{Rigoureux},fonttitle=\bfseries,breakable,#1}
% \newtcolorbox{resultbox}[1][]{colback=RoyalBlue!5,colframe=RoyalBlue!70!black,
%   arc=1mm,title=\textbf{Résultat calculatoire},fonttitle=\bfseries,breakable,#1}
% \newtheorem{theorem}{Théorème}[section]
% \newtheorem{lemma}[theorem]{Lemme}
% \theoremstyle{definition}
% \newtheorem{definition}[theorem]{Définition}
% \begin{document}
% %==============================================================================
% Certificat de primalité inconditionnelle — premier de 29 998 chiffres
% Famille paramétrique  p = 3m(m+1)+1,  m = 2^a·3^b − 1
% Vérification numérique indépendante — mai 2026
%==============================================================================
%
% Utilisation dans le document principal :
%   \input{certificat_29998_fr}
%
% Compilation autonome :
%   pdflatex certificat_29998_fr.tex
%
%==============================================================================

%------------------------------------------------------------------------------
\section{Certificat inconditionnel pour le premier de 29\,998 chiffres}
\label{sec:cert29998}
%------------------------------------------------------------------------------

Cette section présente le certificat de primalité Pocklington–Lehmer
complet et vérifié indépendamment pour le plus grand premier
actuellement produit par \textsc{PrimeQuest}.

%------------------------------------------------------------------------------
\subsection{Paramètres et structure}
%------------------------------------------------------------------------------

\begin{definition}[Instance de 29\,998 chiffres]\label{def:instance29998}
On pose
\[
  a = 19\,435,\qquad b = 19\,173,
\]
et on définit
\[
  m \;=\; 2^{19435}\cdot 3^{19173} - 1
    \qquad\bigl(14\,999 \text{ chiffres décimaux}\bigr),
\]
\[
  p \;=\; 3m(m+1)+1
    \qquad\bigl(29\,998 \text{ chiffres décimaux}\bigr).
\]
\end{definition}

Le rapport $b/a = 19\,173/19\,435 \approx 0{,}9865$ est voisin de l'optimum
théorique $\rho^* = \log_{10}2/\log_{10}3 \approx 0{,}6309$ minimisant le
nombre de chiffres à $a$ fixé ; la valeur retenue ici cible exactement
$29\,998$ chiffres et a été localisée par l'algorithme \textsc{PrimeQuest v4}
(voir Section~\ref{sec:algo}).

\medskip

\noindent\textbf{Identification.}\quad
Les premiers et derniers chiffres décimaux de $p$ sont :
\begin{equation}\label{eq:p29998chiffres}
  p \;=\; 1602087359509060348500037851549319056715\cdots
         16930852697403293697.
\end{equation}

%------------------------------------------------------------------------------
\subsection{Factorisation de $p-1$}
%------------------------------------------------------------------------------

\begin{rigorousbox}
\begin{lemma}[Factorisation de $p-1$]\label{lem:factorisation29998}
Avec $a, b, m, p$ comme dans la Définition~\ref{def:instance29998},
\[
  p - 1 \;=\; \underbrace{2^{19435}\cdot 3^{19174}}_{F}
              \;\cdot\; m,
\]
où $F := 2^{19435}\cdot 3^{19174}$ possède $15\,000$ chiffres décimaux et
vérifie $F > \sqrt{p}$.
\end{lemma}
\end{rigorousbox}

\begin{proof}
Puisque $m + 1 = 2^{19435}\cdot 3^{19173}$,
\[
  p - 1 \;=\; 3m(m+1) \;=\; 3m\cdot 2^{19435}\cdot 3^{19173}
           \;=\; 2^{19435}\cdot 3^{19174}\cdot m.
\]
Pour la borne $F > \sqrt{p}$, on note que $p < 3(m+1)^2$, donc
$\sqrt{p} < \sqrt{3}\,(m+1) < 3(m+1) = 2^{19435}\cdot 3^{19174} = F$.

\noindent Numériquement, en comparant les logarithmes en base $2$ :
\[
  2\log_2 F
  \;=\; 2\bigl(19435 + 19174\log_2 3\bigr)
  \;\approx\; 99\,650{,}14,
\]
\[
  \log_2 p \;\approx\; 99\,648{,}56,
\]
ce qui confirme $F^2 > p$ avec une marge de $\approx 1{,}58$ bit.
\end{proof}

%------------------------------------------------------------------------------
\subsection{Le certificat Pocklington–Lehmer}
%------------------------------------------------------------------------------

Le Théorème~\ref{thm:Pocklington} (rappelé Section~\ref{sec:family})
requiert, pour chaque premier $q \mid F$ — ici $q \in \{2, 3\}$ —, un entier
$w_q$ vérifiant
\[
  w_q^{p-1} \equiv 1 \pmod{p}
  \qquad\text{et}\qquad
  \gcd\!\bigl(w_q^{(p-1)/q} - 1,\; p\bigr) = 1.
\]

\begin{resultbox}[title={\textbf{Certificat} — premier de 29\,998 chiffres}]
\[
  p \;=\; 3m(m+1)+1,\qquad m = 2^{19435}\cdot 3^{19173}-1.
\]

\medskip
\noindent\textbf{Partie connue de $p-1$ :}
$\;F = 2^{19435}\cdot 3^{19174},\quad F > \sqrt{p}.\quad$\checkmark

\medskip
\renewcommand{\arraystretch}{1.4}
\begin{tabular}{@{}c l l@{}}
\toprule
Premier $q$ & Témoin $w_q$ & Conditions vérifiées \\
\midrule
$2$ & $w_2 = 5$ &
  $5^{p-1}\equiv 1\pmod{p}$\quad et\quad
  $\gcd(5^{(p-1)/2}-1,\,p)=1$ \\
$3$ & $w_3 = 7$ &
  $7^{p-1}\equiv 1\pmod{p}$\quad et\quad
  $\gcd(7^{(p-1)/3}-1,\,p)=1$ \\
\bottomrule
\end{tabular}

\medskip
\noindent\textbf{Conclusion :} $p$ est premier (preuve déterministe). \checkmark
\end{resultbox}

%------------------------------------------------------------------------------
\subsection{Vérification numérique indépendante}
%------------------------------------------------------------------------------

Les quatre conditions du certificat ont été vérifiées de façon indépendante
à l'aide de \texttt{Python 3.12} et de son arithmétique entière en précision
arbitraire (backend GMP). Le script ne comporte aucune étape probabiliste ;
toutes les opérations sont des arithmétiques modulaires exactes.

\begin{rigorousbox}[title={\textbf{Attestation de vérification} — calcul indépendant}]
\textbf{Environnement :} Python~3.12,
\texttt{sys.set\_int\_max\_str\_digits(40000)} ;
exponentiations modulaires via \texttt{pow(base, exp, mod)}.\\[4pt]

\textbf{Étape 1 — Reconstruction de $p$.}\quad
Calcul de $m = 2^{19435}\cdot 3^{19173}-1$ (14\,999 chiffres) et de
$p = 3m(m+1)+1$ (29\,998 chiffres) à partir des paramètres bruts.
Premiers 40 chiffres de $p$ :
\texttt{1602087359509060348500037851549319056715} ;
20 derniers : \texttt{16930852697403293697}.
Correspond exactement au certificat stocké.\\[4pt]

\textbf{Étape 2 — Borne Pocklington.}
\[
  2\log_2 F = 99\,650{,}1420 \;>\; \log_2 p = 99\,648{,}5570.\quad\checkmark
\]

\textbf{Étape 3 — Témoin $w_2 = 5$ pour $q = 2$.}
\begin{align*}
  5^{p-1} &\equiv 1 \pmod{p},\quad\checkmark \\
  \gcd(5^{(p-1)/2}-1,\,p) &= 1.\quad\checkmark
\end{align*}

\textbf{Étape 4 — Témoin $w_3 = 7$ pour $q = 3$.}
\begin{align*}
  7^{p-1} &\equiv 1 \pmod{p},\quad\checkmark \\
  \gcd(7^{(p-1)/3}-1,\,p) &= 1.\quad\checkmark
\end{align*}

\textbf{Durée totale :} $6\,945$\,s $\approx 1$\,h\,$56$\,min
(cœur unique, Python GMP, sans parallélisme).\\[4pt]

\textbf{Conclusion :} les quatre conditions Pocklington–Lehmer sont satisfaites.
Le nombre $p$ est \emph{premier de façon inconditionnelle}.
\end{rigorousbox}

%------------------------------------------------------------------------------
\subsection{Comparaison avec le calcul original}
%------------------------------------------------------------------------------

\begin{table}[h]
\centering
\caption{Recherche originale vs.\ vérification indépendante du certificat
         pour le premier de 29\,998 chiffres.}
\label{tab:comparaison29998}
\renewcommand{\arraystretch}{1.3}
\begin{tabular}{@{}lll@{}}
\toprule
 & \textbf{PrimeQuest v4 (recherche)} & \textbf{Vérification indépendante} \\
\midrule
Objectif       & Trouver un candidat premier         & Vérifier le certificat \\
Machine        & ARM64, 9 cœurs (Snapdragon X)       & Cœur unique (Python GMP) \\
Durée          & $11\,146$\,s ($3$\,h\,$06$\,min)   & $6\,945$\,s ($1$\,h\,$56$\,min) \\
Tests MR       & $286$ (pré-filtre probabiliste)     & — (aucune étape probabiliste) \\
Résultat       & Certificat produit                  & Certificat confirmé \\
\bottomrule
\end{tabular}
\end{table}

La vérification est $\approx 1{,}6$ fois plus rapide que la recherche
originale car elle exécute exactement quatre exponentiations modulaires sur
le candidat connu, sans crible, sans filtre du Théorème~2 ni
pré-sélection Miller–Rabin.

%------------------------------------------------------------------------------
\subsection{Rang dans les records}
%------------------------------------------------------------------------------

Le Tableau~\ref{tab:tousresultats} (Section~\ref{sec:resultats}) liste
l'ensemble des quatre certificats \textsc{PrimeQuest}. L'instance de
$29\,998$ chiffres est actuellement le plus grand certificat
Pocklington–Lehmer inconditionnel dans la famille
$p = 3m(m+1)+1$, $m = 2^a\cdot 3^b - 1$,
obtenu sur un ordinateur portable grand public sans calcul distribué ni
matériel spécialisé.

%==============================================================================
% PRÉAMBULE MINIMAL AUTONOME (décommenter pour compilation stand-alone)
%==============================================================================
% \iffalse   % <-- supprimer cette ligne pour activer le mode autonome
% \documentclass[11pt,a4paper]{article}
% \usepackage[utf8]{inputenc}
% \usepackage[T1]{fontenc}
% \usepackage[french]{babel}
% \usepackage{amsmath,amssymb,amsthm,mathtools,booktabs}
% \usepackage{geometry}\geometry{margin=2.5cm}
% \usepackage[dvipsnames]{xcolor}
% \usepackage{tcolorbox}
% \tcbuselibrary{skins,breakable}
% \newtcolorbox{rigorousbox}[1][]{colback=Green!8,colframe=Green!60!black,
%   arc=1mm,title=\textbf{Rigoureux},fonttitle=\bfseries,breakable,#1}
% \newtcolorbox{resultbox}[1][]{colback=RoyalBlue!5,colframe=RoyalBlue!70!black,
%   arc=1mm,title=\textbf{Résultat calculatoire},fonttitle=\bfseries,breakable,#1}
% \newtheorem{theorem}{Théorème}[section]
% \newtheorem{lemma}[theorem]{Lemme}
% \theoremstyle{definition}
% \newtheorem{definition}[theorem]{Définition}
% \begin{document}
% \input{certificat_29998_fr}
% \end{document}
% \fi

% \end{document}
% \fi

% \end{document}
% \fi

% \end{document}
% \fi
